% =====================================================================
% Section 3 -- The RoadScribe system
% Usage: % =====================================================================
% Section 3 -- The RoadScribe system
% Usage: % =====================================================================
% Section 3 -- The RoadScribe system
% Usage: % =====================================================================
% Section 3 -- The RoadScribe system
% Usage: \input{sec3_system} inside an acmart document.
% Preamble requirements (acmart already loads amsmath, amsthm, booktabs, xcolor):
%   \usepackage{algorithm}
%   \usepackage{algpseudocode}
% Bibliography keys used below (add to your .bib):
%   iso2631, iso8855, iso8608, iso13473, newson2009hmm, killick2012pelt,
%   haynes2017crops, kaiser1990tkeo, mohan2008nericell, eriksson2008pothole,
%   menegazzo2021pvs, ggerganov_llamacpp
% =====================================================================
\providecommand{\R}{\mathbb{R}}
\providecommand{\ind}{\mathbf{1}}
\providecommand{\TBD}[1][]{\textcolor{red}{[TBD#1]}}

\section{The RoadScribe System}
\label{sec:system}

RoadScribe turns a trip recorded by a commodity smartphone into a short narrative and answers follow-up questions about it.
The pipeline has three layers (Figure~\ref{fig:arch}).
Layer~1 converts inertial and GNSS samples into vehicle-frame ride measures and candidate events.
Layer~2 anchors these measures to the road network and partitions the trip into homogeneous road narrative units (HRNUs).
Layer~3 compiles the units into a token-budgeted capsule and a machine-checkable claim ledger, and generates text with an on-device small language model (SLM) under a plan--ground--realize protocol.

Four principles shape the design.
(P1)~Every number, place name and event in the output can be traced to a ledger fact.
(P2)~The language model decides what to say and how to say it, but never supplies a value; values come from the ledger.
(P3)~Uncertainty estimated in Layers~1--2 reaches the text as calibrated hedges.
(P4)~When verification fails, a deterministic realizer takes over, so the SLM can add coverage but cannot remove faithfulness.

\begin{figure}[t]
  \centering
  \fbox{\parbox{0.95\linewidth}{\centering\vspace{2em}Architecture figure placeholder: IMU+GNSS $\rightarrow$ Layer 1 (orientation, ISO weighting, events) $\rightarrow$ Layer 2 (map matching, MI-PELT, HRNUs) $\rightarrow$ Layer 3 (ledger, KST capsule, plan--ground--realize, verifier, template fallback) $\rightarrow$ narrative / answers\vspace{2em}}}
  \caption{RoadScribe architecture.}
  \label{fig:arch}
\end{figure}

% ---------------------------------------------------------------------
\subsection{Layer 1: Kinetic Signal Processing}
\label{sec:layer1}

\paragraph{Notation.}
The IMU delivers specific force $\mathbf f^b[n]\in\R^3$ and angular rate $\boldsymbol\omega^b[n]\in\R^3$ in the phone body frame $b$ at rate $f_s$ ($T_s=1/f_s$).
GNSS delivers speed $v_k$ at times $t_k$ (1~Hz).
The vehicle frame $v$ follows ISO~8855~\cite{iso8855} ($x$ forward, $y$ left, $z$ up).
We estimate $\mathbf R_{vb}\in SO(3)$ and the gravity-free acceleration
\begin{equation}
  \mathbf a^v[n]=\mathbf R_{vb}\,\mathbf f^b[n]-g\,\mathbf e_z .
  \label{eq:av}
\end{equation}
Speed is interpolated to IMU time with a monotone cubic (PCHIP), giving $v(t)$ and $a^G_{\mathrm{lon}}(t)=\dot v(t)$.

\paragraph{Mount episodes.}
A phone in a holder is static relative to the vehicle except when handled.
We mark candidate handling where the gyroscope reports rotation that the vehicle does not perform, $\lVert\bar{\boldsymbol\omega}^b[n]\rVert-|\dot\psi^G(t_n)|>\Omega_h$, with $\dot\psi^G$ the GNSS course rate.
A candidate becomes an episode boundary when the quasi-static gravity directions on its two sides differ by more than $\theta_h$.
Orientation is estimated per episode $e$.

\paragraph{Gravity and tilt.}
Let $\bar{\mathbf f}^b$ and $\bar{\boldsymbol\omega}^b$ be low-passed at 0.5~Hz. The quasi-static set of episode $e$ is
\begin{equation}
\mathcal Q_e=\bigl\{n:\ \bigl|\lVert\bar{\mathbf f}^b[n]\rVert-g\bigr|<\epsilon_g,\ \lVert\bar{\boldsymbol\omega}^b[n]\rVert<\epsilon_\omega,\ |a^G_{\mathrm{lon}}(t_n)|<\epsilon_a,\ v(t_n)>v_{\min}\bigr\},
\end{equation}
and the up direction is $\mathbf u_e=\hat{\mathbf g}_e/\lVert\hat{\mathbf g}_e\rVert$ with $\hat{\mathbf g}_e$ the geometric median of $\{\bar{\mathbf f}^b[n]:n\in\mathcal Q_e\}$, which is robust to outliers and equivariant to rotation.
Gating matters: an ungated low-pass estimate absorbs sustained acceleration and tilts by $\arctan(a/g)$, i.e.\ $17^\circ$ at $3~\mathrm{m\,s^{-2}}$.
With $\mathbf k=\mathbf u_e\times\mathbf e_z/\lVert\mathbf u_e\times\mathbf e_z\rVert$ and $\theta=\operatorname{atan2}(\lVert\mathbf u_e\times\mathbf e_z\rVert,\mathbf u_e^\top\mathbf e_z)$, Rodrigues' formula gives the leveling rotation
\begin{equation}
\mathbf R_1=\mathbf I+\sin\theta\,[\mathbf k]_\times+(1-\cos\theta)\,[\mathbf k]_\times^2 .
\end{equation}
The route-average road grade biases this tilt (a 5\% grade is $2.9^\circ$).
The resulting error is at most 0.13\% on the vertical channel, but up to $0.5~\mathrm{m\,s^{-2}}$ on the longitudinal channel.
We therefore take longitudinal maneuver thresholds from GNSS rather than from the IMU.

\paragraph{Heading.}
After leveling, the horizontal axes differ from the vehicle axes by an unknown yaw $\psi$.
The leveled yaw rate $\omega^\ell_z=(\mathbf R_1\bar{\boldsymbol\omega}^b)_z$ does not depend on $\psi$.
This yields two GNSS-referenced targets at every fix: the longitudinal acceleration $a^G_{\mathrm{lon},k}$ and the centripetal acceleration $a^G_{\mathrm{lat},k}=v_k\,\omega^\ell_z(t_k)$.
With the leveled horizontal specific force $\mathbf h_k\in\R^2$ and $\mathbf r_k=(a^G_{\mathrm{lon},k},a^G_{\mathrm{lat},k})^\top$,
\begin{equation}
\hat\psi=\arg\min_{\psi}\sum_k\bigl\lVert\mathbf R_2(\psi)\mathbf h_k-\mathbf r_k\bigr\rVert^2
=\operatorname{atan2}\bigl(M_{21}-M_{12},\ M_{11}+M_{22}\bigr),\qquad \mathbf M=\sum_k\mathbf h_k\mathbf r_k^\top,
\label{eq:yaw}
\end{equation}
where $\mathbf R_2(\psi)=\bigl[\begin{smallmatrix}\cos\psi&\sin\psi\\-\sin\psi&\cos\psi\end{smallmatrix}\bigr]$.
The closed form follows from maximizing $\operatorname{tr}(\mathbf R_2\mathbf M)$.
Only fixes with $v_k>3~\mathrm{m\,s^{-1}}$ are used.
Finally $\mathbf R_{vb}=\mathbf R_z(\hat\psi)\mathbf R_1$, and a bootstrap over fixes provides the orientation confidence passed to Layer~3.

\paragraph{Speed decoupling.}
Vibration grows with speed.
For a road whose elevation PSD follows ISO~8608~\cite{iso8608}, $G_d(n)=G_d(n_0)(n/n_0)^{-w}$, the temporal input PSD at speed $v$ scales as $v^{w-1}f^{-w}$.
For a fixed vehicle response the RMS response therefore scales as $v^{(w-1)/2}$.
Rather than assume $w$, we estimate the elasticity with the window-level model $\log a_{w,i}=\alpha_{j(i)}+\gamma\log v_i+\varepsilon_i$, where $j$ indexes homogeneous segments, and report
\begin{equation}
\tilde a_w=a_w\,(v_{\mathrm{ref}}/\bar v)^{\hat\gamma},\qquad \bar v\ge v_{\min},\ v_{\mathrm{ref}}=10~\mathrm{m\,s^{-1}}.
\end{equation}
The raw $a_w$ describes how the ride felt and $\tilde a_w$ describes the road; narratives keep the two apart.

\paragraph{Frequency weighting.}
ISO~2631-1~\cite{iso2631} weights seated whole-body vibration with $W_k$ (vertical) and $W_d$ (horizontal).
Each is a product of band-limiting filters, an acceleration--velocity transition and, for $W_k$, an upward step ($p=j2\pi f$, $\omega_i=2\pi f_i$, $Q_1=1/\sqrt2$):
\begin{align}
H_h(p)&=\frac{p^2}{p^2+\omega_1p/Q_1+\omega_1^2}, &
H_l(p)&=\frac{\omega_2^2}{p^2+\omega_2p/Q_1+\omega_2^2}, \nonumber\\
H_t(p)&=\frac{(\omega_4^2/\omega_3)(p+\omega_3)}{p^2+\omega_4p/Q_4+\omega_4^2}, &
H_s(p)&=\frac{p^2+\omega_5p/Q_5+\omega_5^2}{p^2+\omega_6p/Q_6+\omega_6^2},
\end{align}
with $W=H_hH_lH_tH_s$ and the parameters of Table~\ref{tab:iso}.
Narratives are produced after the trip, so we apply $W$ in the frequency domain, $a_w=\mathcal F^{-1}\{\mathcal F\{a\}\,W(j2\pi f)\}$ with 10~s of zero padding.
This reproduces the analog magnitude and phase exactly below $f_s/2$.
For streaming use, a bilinear transform per second-order section, pre-warped at each section's corner frequency, is an alternative.
At $f_s=100$~Hz its magnitude error is at most 0.72~dB ($W_k$) and 1.22~dB ($W_d$) up to 20~Hz, but it reaches 7--8~dB at 40~Hz and underestimates impulsive peaks by 3.8\%.
At $f_s=400$~Hz the error stays within 0.23~dB up to 40~Hz.
The 100~Hz band limit $H_l$ cannot be realized for $f_s\le200$~Hz; omitting it changes $W$ by less than 0.11~dB below 40~Hz.
At $f_s=100$~Hz our analysis band ends near 40~Hz, below the 80~Hz upper limit of the standard.

\begin{table}[t]
\caption{ISO 2631-1 frequency-weighting parameters (Annex A). Frequencies in Hz.}
\label{tab:iso}
\begin{tabular}{lccccccccc}
\toprule
 & $f_1$ & $f_2$ & $f_3$ & $f_4$ & $Q_4$ & $f_5$ & $Q_5$ & $f_6$ & $Q_6$\\
\midrule
$W_k$ & 0.4 & 100 & 12.5 & 12.5 & 0.63 & 2.37 & 0.91 & 3.35 & 0.91\\
$W_d$ & 0.4 & 100 & 2.0 & 2.0 & 0.63 & -- & -- & -- & --\\
\bottomrule
\end{tabular}
\end{table}

\paragraph{Vibration measures.}
For each HRNU of duration $T$ and axis $j$ we compute
\begin{gather}
a_{w,j}=\Bigl(\tfrac1T\textstyle\int_0^Ta_{w,j}^2\,dt\Bigr)^{1/2},\qquad
a_v=\bigl(a_{w,x}^2+a_{w,y}^2+a_{w,z}^2\bigr)^{1/2},\\
\mathrm{MTVV}_j=\max_{t_0}\Bigl(\tfrac1\tau\textstyle\int_{t_0-\tau}^{t_0}a_{w,j}^2\,dt\Bigr)^{1/2}\ (\tau=1~\mathrm{s}),\qquad
\mathrm{VDV}_j=\Bigl(\textstyle\int_0^Ta_{w,j}^4\,dt\Bigr)^{1/4},\qquad
\mathrm{CF}_j=\frac{\max_t|a_{w,j}(t)|}{a_{w,j}}.
\end{gather}
Following the standard's evaluation rule, $a_w$ is always reported.
MTVV and VDV are added, and the unit is flagged as shock-dominated, when
\begin{equation}
S_j=\ind\bigl[\mathrm{CF}_j>9\ \lor\ \mathrm{MTVV}_j/a_{w,j}>1.5\ \lor\ \mathrm{VDV}_j/(a_{w,j}T^{1/4})>1.75\bigr].
\end{equation}
For a zero-mean stationary signal $\mathrm{VDV}/(a_wT^{1/4})=\kappa^{1/4}$, where $\kappa$ is the kurtosis.
The VDV criterion is therefore the condition $\kappa>9.38$; Gaussian vibration gives a ratio of 1.32.
Comfort descriptors follow the informative comfort scale of the standard (Annex~C).
Its bands overlap, so a single label uses geometric-mean cut-points $\boldsymbol\theta=(0.315,0.561,0.894,1.414,2.236)~\mathrm{m\,s^{-2}}$.
ISO~2631-1 defines comfort at the seat--occupant interface, whereas a phone in a holder measures the cabin structure.
We therefore call these quantities \emph{weighted cabin acceleration} and validate them against a phone fixed at the seat rail (Section~\ref{sec:data}).
Below about 1~Hz, $a_{w,x}$ and $a_{w,y}$ mostly reflect braking and cornering, and we report them as driving rather than road measures.

\paragraph{Modal band ratios.}
From the Welch PSD of $a^v_z$ (Hann, 5.12~s segments, 50\% overlap) we compute energies in a body band $[0.8,2.5]$~Hz, a wheel-hop band centred on the vehicle's wheel-hop peak $\hat f_{wh}\in[6,20]$~Hz, and a high band $[16,40]$~Hz.
From these we form $R_{wb}=\log(E_{\mathrm{wheel}}/E_{\mathrm{body}})$ and $R_{hb}=\log(E_{\mathrm{high}}/E_{\mathrm{body}})$.
Through $\lambda=v/f$ these bands correspond to wavelengths from about 0.25~m to 12.5~m at 10~m\,s$^{-1}$, i.e.\ unevenness and megatexture, not pavement macrotexture (0.5--50~mm~\cite{iso13473}).
They describe how the vehicle's modes are excited, and we use them within a vehicle and for segmentation.

\paragraph{Impulsive events and maneuver context.}
We band-pass the vertical channel, $\tilde z=\mathrm{BP}_{[4,25]\,\mathrm{Hz}}(a^v_z)$, and apply the Teager--Kaiser energy operator~\cite{kaiser1990tkeo},
\begin{equation}
\Psi[n]=\tilde z[n]^2-\tilde z[n-1]\,\tilde z[n+1],\qquad E[n]=\bigl(h_{50\,\mathrm{ms}}*\max(\Psi,0)\bigr)[n].
\end{equation}
For $\tilde z=A\cos(\Omega n+\phi)$, $\Psi=A^2\sin^2\Omega$, so the operator responds to both amplitude and frequency.
Candidates are maximal runs where $E$ exceeds $\theta[n]=\operatorname{med}_W(E)+\kappa\cdot1.4826\operatorname{MAD}_W(E)$ over a 10~s centred window.
The adaptive threshold follows the background, which is much higher on cobblestone than on asphalt.
Runs closer than 0.3~s are merged, and runs longer than $D_{\max}=0.5~\mathrm{s}+3~\mathrm{m}/\bar v$ are passed to segmentation as rough stretches.
The maneuver state $m_e$ is one of \textsc{brake}, \textsc{accel}, \textsc{turn}, \textsc{swerve}, \textsc{stop-go} or \textsc{cruise}.
Longitudinal states come from GNSS ($a^G_{\mathrm{lon}}$) and lateral ones from the yaw rate ($v\,\omega^v_z$); \textsc{stop-go} requires $v<2~\mathrm{m\,s^{-1}}$ and a longitudinal jerk $|j_x|>J_{sg}$.
With $\rho_e=E_{[8,25]}(e)/E_{[0.5,4]}(e)$ the band-energy ratio of $a^v_z$ in the event window, each candidate is labelled
\begin{equation}
\ell(e)=
\begin{cases}
\text{\textsc{artifact}} & m_e\in\{\text{\textsc{stop-go}},\text{\textsc{accel}}\}\ \land\ \max_{n\in e}|j_x[n]|>J_x\ \land\ \rho_e<\rho_0,\\
\text{\textsc{defect}} & \lnot\,\text{\textsc{artifact}}\ \land\ \rho_e\ge\rho_0\ \land\ d_e\le D_{\max},\\
\text{\textsc{ambiguous}} & \text{otherwise}.
\end{cases}
\label{eq:gate}
\end{equation}
Braking and turning never veto a defect.
In our data they do not increase vertical vibration (Section~\ref{sec:l1eval}), and drivers slow down \emph{before} speed bumps, so a braking veto would remove true events.
Acceleration is gated because it is the only maneuver that adds a dashboard-only high-band component.
A swerve without a vertical candidate within $\pm2$~s yields a separate low-confidence \emph{avoidance} event.
A logistic model trained on held-out routes maps the event features to $p_{\mathrm{defect}}$.

% ---------------------------------------------------------------------
\subsection{Layer 2: Spatial Grounding and Segmentation}
\label{sec:layer2}

\paragraph{Map matching.}
We match GNSS fixes to OpenStreetMap with a hidden Markov model~\cite{newson2009hmm}.
The states are candidate projections $x_{k,i}$ within 50~m of fix $z_k$, the emissions are $\mathcal N(\lVert z_k-x_{k,i}\rVert_{gc};0,\sigma_z^2)$, and the transitions are $\beta^{-1}\exp(-|\lVert z_k-z_{k+1}\rVert_{gc}-\lVert x_{k,i}-x_{k+1,j}\rVert_{\mathrm{route}}|/\beta)$, decoded with Viterbi.
Forward--backward marginals give a per-fix match confidence.
The matched arc length $s(t)$ replaces $\int v\,dt$.
Each matched way contributes its \texttt{surface}, \texttt{smoothness}, \texttt{highway} and \texttt{name} tags, and nodes such as \texttt{traffic\_calming} and level crossings become candidate anchors.

\paragraph{Map-informed change-point detection.}
We resample Layer-1 features on a 5~m along-track grid, $\mathbf y_j=[\log\tilde a_{w,z},\log a_{w,z},R_{wb},R_{hb},\text{event rate},\log a_{w,y}]_j$, robust-standardized per trip.
Speed is excluded so that boundaries do not follow speed changes.
Segmentation solves
\begin{equation}
\min_{m,\,\tau_{1:m}}\ \sum_{i=1}^{m+1}C\bigl(\mathbf y_{\tau_{i-1}+1:\tau_i}\bigr)+\sum_{i=1}^{m}\beta(\tau_i),\qquad
C(\mathbf y_{a:b})=\sum_{j=a}^{b}\lVert\mathbf y_j-\bar{\mathbf y}_{a:b}\rVert^2_{\mathbf D^{-1}},
\label{eq:pelt}
\end{equation}
with a penalty that falls near OpenStreetMap transitions $b\in\mathcal B$ (surface-tag changes, road-class changes, junctions):
\begin{equation}
\beta(\tau)=\beta_0\bigl(1-\lambda\,\pi(\tau)\bigr),\qquad
\pi(\tau)=\max_{b\in\mathcal B}w_b\exp\!\Bigl(-\frac{(s_\tau-s_b)^2}{2\sigma_b^2}\Bigr),\qquad \lambda\in[0,1).
\end{equation}
Setting $\lambda=0$ recovers PELT~\cite{killick2012pelt}.
Let $F(0)=0$, $G(t)=F(t)+\beta(t)\ind[t>0]$ and $F(s)=\min_{t\in\mathcal R_s}\{G(t)+C(\mathbf y_{t+1:s})\}$.

\begin{lemma}[Pruning under location-dependent penalties]
\label{lem:prune}
Suppose $C(\mathbf y_{t+1:s})+C(\mathbf y_{s+1:T})\le C(\mathbf y_{t+1:T})$ for all $t<s<T$.
If $G(t)+C(\mathbf y_{t+1:s})\ge G(s)$, then $t$ is not the optimal last change point for any $T>s$ and can be removed from $\mathcal R$.
\end{lemma}
\begin{proof}
For $T>s$, $G(t)+C(\mathbf y_{t+1:T})\ge G(t)+C(\mathbf y_{t+1:s})+C(\mathbf y_{s+1:T})\ge G(s)+C(\mathbf y_{s+1:T})$.
\end{proof}

The usual condition $F(t)+C(\mathbf y_{t+1:s})\ge F(s)$ coincides with Lemma~\ref{lem:prune} only for a constant penalty.
With a lower penalty at map transitions it discards valid candidates.
In 300 random series it returned a sub-optimal partition every time, whereas the rule of Lemma~\ref{lem:prune} matched exhaustive optimal partitioning in all 300.
We choose $\beta_0=c\,d\ln n$ and $(c,\lambda)$ by leave-one-route-out validation, and report the penalty path~\cite{haynes2017crops}.

\paragraph{Arbitration and HRNUs.}
A change point within 25~m of a relevant map transition is snapped to it: the map fixes the location and the signal decides existence.
A map transition without a change point is kept only if a two-sample Hotelling test on $\mathbf y$ across it rejects at $\alpha=0.01$; otherwise every way split in the map would start a new unit.
Units shorter than 50~m are merged into their closer neighbour.
Each boundary is typed as map+signal, signal-only or map-only.
Each unit carries a map--signal disagreement score $1-P(\text{surface}=\text{OSM tag}\mid\mathbf y)$ when the map has a surface tag.
The resulting HRNUs, with their measures, events, anchors and confidences, are the input to Layer~3.

% ---------------------------------------------------------------------
\subsection{Layer 3: Kinetic--Spatial Tokenization and Grounded Generation}
\label{sec:layer3}

\paragraph{Claim ledger.}
Layer~3 first converts HRNUs and events into a ledger $\mathcal L$ of typed facts with identifiers, values, tolerances and confidences.
The types are severity, numeric, event, anchor, comparison, missing and metadata.
A \textsc{missing} fact records that a value is unavailable and why, e.g.\ a 48~s GNSS gap.
The ledger is the reference for both the run-time verifier and the evaluation.

\paragraph{Kinetic--Spatial Tokenizer.}
Algorithm~\ref{alg:kst} compiles the ledger into a capsule that fits a token budget $B$.
Each unit has a salience $\sigma_m$ and three serializations (omit, brief, full) with token costs $c_{m,\ell}$ measured with the target tokenizer.
Selection solves a multiple-choice knapsack exactly by dynamic programming:
\begin{equation}
V_m(b)=\max_{\ell\in\Lambda_m}\bigl\{V_{m-1}(b-c_{m,\ell})+\sigma_m\eta_\ell\bigr\},\qquad V_0(b)=0,
\label{eq:knap}
\end{equation}
where $\Lambda_m$ excludes \textsc{omit} for units in the must-mention set $\mathcal M$ and $\eta_{\textsc{omit}}=0$.
The cost is $O(3MB)$ for $M$ units.

\begin{algorithm}[t]
\caption{Kinetic--Spatial Tokenizer (KST)}
\label{alg:kst}
\begin{algorithmic}[1]
\Require HRNUs $\mathcal U=(u_1,\dots,u_M)$ in route order; events $\mathcal E$; anchors $\mathcal A$; optional history $\mathcal H$; audience $\pi$; budget $B$; tokenizer $\mathrm{tok}$
\Ensure capsule $\kappa$ with $|\mathrm{tok}(\kappa)|\le B$; ledger $\mathcal L$
\State $\mathcal L\gets\textsc{MetaFacts}(\text{channels},\text{map coverage},\pi)$
\For{$m=1,\dots,M$}
  \State $k_m\gets\textsc{IsoClass}(a_{v,m})$;\quad $\tilde k_m\gets\textsc{IsoClass}(\tilde a_{w,m})$ \Comment{felt vs.\ road}
  \State $\mathcal E_m\gets\{e\in\mathcal E: s(e)\in[s^0_m,s^1_m)\}$
  \State $\alpha_m\gets$ nearest anchor in $\mathcal A$ within $r_{\max}$ of $u_m$, or $\varnothing$
  \State $\sigma_m\gets w_1k_m+w_2\sum_{e\in\mathcal E_m}p_e\,\mathrm{sev}(e)+w_3\,\Delta_{\mathcal H}(u_m)+w_4\,\delta^{\mathrm{map}}_m$
  \State $\mathcal L\gets\mathcal L\cup\textsc{Facts}(u_m,k_m,\tilde k_m,\mathcal E_m,\alpha_m)$ \Comment{including \textsc{missing} facts}
  \For{$\ell\in\{\textsc{omit},\textsc{brief},\textsc{full}\}$}
    \State $c_{m,\ell}\gets|\mathrm{tok}(\textsc{Ser}_\ell(u_m))|$
  \EndFor
\EndFor
\State $\mathcal M\gets\{m:\exists e\in\mathcal E_m,\ p_e\,\mathrm{sev}(e)\ge\eta^\star\}$ \Comment{must-mention units}
\State $B'\gets B-|\mathrm{tok}(\textsc{Header})|-c_{\mathrm{rest}}$
\State $\mathbf x^\star\gets\textsc{McKnapsack}(\{c_{m,\ell}\},\{\sigma_m\eta_\ell\},B',\mathcal M)$ \Comment{Eq.~\eqref{eq:knap}}
\If{$\mathbf x^\star$ is infeasible}
  \State merge adjacent units outside $\mathcal M$ and repeat; if still infeasible, increase $B$ and log it
\EndIf
\State $\kappa\gets\textsc{Header}\ \|\ \textsc{Ser}_{x^\star_1}(u_1)\ \|\cdots\|\ \textsc{Ser}_{x^\star_M}(u_M)\ \|\ \textsc{Rest}(\{u_m:x^\star_m=\textsc{omit}\})$
\State \Return $\kappa,\ \mathcal L$
\end{algorithmic}
\end{algorithm}

\paragraph{Plan--ground--realize.}
Generation runs in four steps.
\emph{Plan:} the SLM emits a JSON plan listing which units and facts to mention, their order, a rhetorical relation for each (elaboration, contrast, cause, comparison with history) and a hedge.
Decoding is constrained by a grammar generated per capsule from the ledger identifiers~\cite{ggerganov_llamacpp}, so the plan is always well-formed and cannot reference a unit or event that does not exist.
\emph{Ground:} deterministic code expands each plan item into a fact frame with values, names and units taken from $\mathcal L$.
\emph{Realize:} the SLM writes text from the fact frames, either with placeholders for numbers and names that are filled afterwards, or as free text.
\emph{Verify:} claims extracted from the text (Section~\ref{sec:metrics}) are checked against $\mathcal L$.
A failing sentence is regenerated at most twice and then replaced by the template realization of the same plan item.
The SLM therefore controls selection, ordering, aggregation and phrasing, but not values.

\paragraph{Hedging.}
The plan assigns hedges from fact confidence $p$: assertive for $p\ge0.9$, \emph{likely} for $0.7\le p<0.9$, \emph{possibly} for $0.4\le p<0.7$; facts with $p<0.4$ are mentioned only when asked about.
The verifier rejects a sentence whose hedge is stronger than its confidence band allows.
 inside an acmart document.
% Preamble requirements (acmart already loads amsmath, amsthm, booktabs, xcolor):
%   \usepackage{algorithm}
%   \usepackage{algpseudocode}
% Bibliography keys used below (add to your .bib):
%   iso2631, iso8855, iso8608, iso13473, newson2009hmm, killick2012pelt,
%   haynes2017crops, kaiser1990tkeo, mohan2008nericell, eriksson2008pothole,
%   menegazzo2021pvs, ggerganov_llamacpp
% =====================================================================
\providecommand{\R}{\mathbb{R}}
\providecommand{\ind}{\mathbf{1}}
\providecommand{\TBD}[1][]{\textcolor{red}{[TBD#1]}}

\section{The RoadScribe System}
\label{sec:system}

RoadScribe turns a trip recorded by a commodity smartphone into a short narrative and answers follow-up questions about it.
The pipeline has three layers (Figure~\ref{fig:arch}).
Layer~1 converts inertial and GNSS samples into vehicle-frame ride measures and candidate events.
Layer~2 anchors these measures to the road network and partitions the trip into homogeneous road narrative units (HRNUs).
Layer~3 compiles the units into a token-budgeted capsule and a machine-checkable claim ledger, and generates text with an on-device small language model (SLM) under a plan--ground--realize protocol.

Four principles shape the design.
(P1)~Every number, place name and event in the output can be traced to a ledger fact.
(P2)~The language model decides what to say and how to say it, but never supplies a value; values come from the ledger.
(P3)~Uncertainty estimated in Layers~1--2 reaches the text as calibrated hedges.
(P4)~When verification fails, a deterministic realizer takes over, so the SLM can add coverage but cannot remove faithfulness.

\begin{figure}[t]
  \centering
  \fbox{\parbox{0.95\linewidth}{\centering\vspace{2em}Architecture figure placeholder: IMU+GNSS $\rightarrow$ Layer 1 (orientation, ISO weighting, events) $\rightarrow$ Layer 2 (map matching, MI-PELT, HRNUs) $\rightarrow$ Layer 3 (ledger, KST capsule, plan--ground--realize, verifier, template fallback) $\rightarrow$ narrative / answers\vspace{2em}}}
  \caption{RoadScribe architecture.}
  \label{fig:arch}
\end{figure}

% ---------------------------------------------------------------------
\subsection{Layer 1: Kinetic Signal Processing}
\label{sec:layer1}

\paragraph{Notation.}
The IMU delivers specific force $\mathbf f^b[n]\in\R^3$ and angular rate $\boldsymbol\omega^b[n]\in\R^3$ in the phone body frame $b$ at rate $f_s$ ($T_s=1/f_s$).
GNSS delivers speed $v_k$ at times $t_k$ (1~Hz).
The vehicle frame $v$ follows ISO~8855~\cite{iso8855} ($x$ forward, $y$ left, $z$ up).
We estimate $\mathbf R_{vb}\in SO(3)$ and the gravity-free acceleration
\begin{equation}
  \mathbf a^v[n]=\mathbf R_{vb}\,\mathbf f^b[n]-g\,\mathbf e_z .
  \label{eq:av}
\end{equation}
Speed is interpolated to IMU time with a monotone cubic (PCHIP), giving $v(t)$ and $a^G_{\mathrm{lon}}(t)=\dot v(t)$.

\paragraph{Mount episodes.}
A phone in a holder is static relative to the vehicle except when handled.
We mark candidate handling where the gyroscope reports rotation that the vehicle does not perform, $\lVert\bar{\boldsymbol\omega}^b[n]\rVert-|\dot\psi^G(t_n)|>\Omega_h$, with $\dot\psi^G$ the GNSS course rate.
A candidate becomes an episode boundary when the quasi-static gravity directions on its two sides differ by more than $\theta_h$.
Orientation is estimated per episode $e$.

\paragraph{Gravity and tilt.}
Let $\bar{\mathbf f}^b$ and $\bar{\boldsymbol\omega}^b$ be low-passed at 0.5~Hz. The quasi-static set of episode $e$ is
\begin{equation}
\mathcal Q_e=\bigl\{n:\ \bigl|\lVert\bar{\mathbf f}^b[n]\rVert-g\bigr|<\epsilon_g,\ \lVert\bar{\boldsymbol\omega}^b[n]\rVert<\epsilon_\omega,\ |a^G_{\mathrm{lon}}(t_n)|<\epsilon_a,\ v(t_n)>v_{\min}\bigr\},
\end{equation}
and the up direction is $\mathbf u_e=\hat{\mathbf g}_e/\lVert\hat{\mathbf g}_e\rVert$ with $\hat{\mathbf g}_e$ the geometric median of $\{\bar{\mathbf f}^b[n]:n\in\mathcal Q_e\}$, which is robust to outliers and equivariant to rotation.
Gating matters: an ungated low-pass estimate absorbs sustained acceleration and tilts by $\arctan(a/g)$, i.e.\ $17^\circ$ at $3~\mathrm{m\,s^{-2}}$.
With $\mathbf k=\mathbf u_e\times\mathbf e_z/\lVert\mathbf u_e\times\mathbf e_z\rVert$ and $\theta=\operatorname{atan2}(\lVert\mathbf u_e\times\mathbf e_z\rVert,\mathbf u_e^\top\mathbf e_z)$, Rodrigues' formula gives the leveling rotation
\begin{equation}
\mathbf R_1=\mathbf I+\sin\theta\,[\mathbf k]_\times+(1-\cos\theta)\,[\mathbf k]_\times^2 .
\end{equation}
The route-average road grade biases this tilt (a 5\% grade is $2.9^\circ$).
The resulting error is at most 0.13\% on the vertical channel, but up to $0.5~\mathrm{m\,s^{-2}}$ on the longitudinal channel.
We therefore take longitudinal maneuver thresholds from GNSS rather than from the IMU.

\paragraph{Heading.}
After leveling, the horizontal axes differ from the vehicle axes by an unknown yaw $\psi$.
The leveled yaw rate $\omega^\ell_z=(\mathbf R_1\bar{\boldsymbol\omega}^b)_z$ does not depend on $\psi$.
This yields two GNSS-referenced targets at every fix: the longitudinal acceleration $a^G_{\mathrm{lon},k}$ and the centripetal acceleration $a^G_{\mathrm{lat},k}=v_k\,\omega^\ell_z(t_k)$.
With the leveled horizontal specific force $\mathbf h_k\in\R^2$ and $\mathbf r_k=(a^G_{\mathrm{lon},k},a^G_{\mathrm{lat},k})^\top$,
\begin{equation}
\hat\psi=\arg\min_{\psi}\sum_k\bigl\lVert\mathbf R_2(\psi)\mathbf h_k-\mathbf r_k\bigr\rVert^2
=\operatorname{atan2}\bigl(M_{21}-M_{12},\ M_{11}+M_{22}\bigr),\qquad \mathbf M=\sum_k\mathbf h_k\mathbf r_k^\top,
\label{eq:yaw}
\end{equation}
where $\mathbf R_2(\psi)=\bigl[\begin{smallmatrix}\cos\psi&\sin\psi\\-\sin\psi&\cos\psi\end{smallmatrix}\bigr]$.
The closed form follows from maximizing $\operatorname{tr}(\mathbf R_2\mathbf M)$.
Only fixes with $v_k>3~\mathrm{m\,s^{-1}}$ are used.
Finally $\mathbf R_{vb}=\mathbf R_z(\hat\psi)\mathbf R_1$, and a bootstrap over fixes provides the orientation confidence passed to Layer~3.

\paragraph{Speed decoupling.}
Vibration grows with speed.
For a road whose elevation PSD follows ISO~8608~\cite{iso8608}, $G_d(n)=G_d(n_0)(n/n_0)^{-w}$, the temporal input PSD at speed $v$ scales as $v^{w-1}f^{-w}$.
For a fixed vehicle response the RMS response therefore scales as $v^{(w-1)/2}$.
Rather than assume $w$, we estimate the elasticity with the window-level model $\log a_{w,i}=\alpha_{j(i)}+\gamma\log v_i+\varepsilon_i$, where $j$ indexes homogeneous segments, and report
\begin{equation}
\tilde a_w=a_w\,(v_{\mathrm{ref}}/\bar v)^{\hat\gamma},\qquad \bar v\ge v_{\min},\ v_{\mathrm{ref}}=10~\mathrm{m\,s^{-1}}.
\end{equation}
The raw $a_w$ describes how the ride felt and $\tilde a_w$ describes the road; narratives keep the two apart.

\paragraph{Frequency weighting.}
ISO~2631-1~\cite{iso2631} weights seated whole-body vibration with $W_k$ (vertical) and $W_d$ (horizontal).
Each is a product of band-limiting filters, an acceleration--velocity transition and, for $W_k$, an upward step ($p=j2\pi f$, $\omega_i=2\pi f_i$, $Q_1=1/\sqrt2$):
\begin{align}
H_h(p)&=\frac{p^2}{p^2+\omega_1p/Q_1+\omega_1^2}, &
H_l(p)&=\frac{\omega_2^2}{p^2+\omega_2p/Q_1+\omega_2^2}, \nonumber\\
H_t(p)&=\frac{(\omega_4^2/\omega_3)(p+\omega_3)}{p^2+\omega_4p/Q_4+\omega_4^2}, &
H_s(p)&=\frac{p^2+\omega_5p/Q_5+\omega_5^2}{p^2+\omega_6p/Q_6+\omega_6^2},
\end{align}
with $W=H_hH_lH_tH_s$ and the parameters of Table~\ref{tab:iso}.
Narratives are produced after the trip, so we apply $W$ in the frequency domain, $a_w=\mathcal F^{-1}\{\mathcal F\{a\}\,W(j2\pi f)\}$ with 10~s of zero padding.
This reproduces the analog magnitude and phase exactly below $f_s/2$.
For streaming use, a bilinear transform per second-order section, pre-warped at each section's corner frequency, is an alternative.
At $f_s=100$~Hz its magnitude error is at most 0.72~dB ($W_k$) and 1.22~dB ($W_d$) up to 20~Hz, but it reaches 7--8~dB at 40~Hz and underestimates impulsive peaks by 3.8\%.
At $f_s=400$~Hz the error stays within 0.23~dB up to 40~Hz.
The 100~Hz band limit $H_l$ cannot be realized for $f_s\le200$~Hz; omitting it changes $W$ by less than 0.11~dB below 40~Hz.
At $f_s=100$~Hz our analysis band ends near 40~Hz, below the 80~Hz upper limit of the standard.

\begin{table}[t]
\caption{ISO 2631-1 frequency-weighting parameters (Annex A). Frequencies in Hz.}
\label{tab:iso}
\begin{tabular}{lccccccccc}
\toprule
 & $f_1$ & $f_2$ & $f_3$ & $f_4$ & $Q_4$ & $f_5$ & $Q_5$ & $f_6$ & $Q_6$\\
\midrule
$W_k$ & 0.4 & 100 & 12.5 & 12.5 & 0.63 & 2.37 & 0.91 & 3.35 & 0.91\\
$W_d$ & 0.4 & 100 & 2.0 & 2.0 & 0.63 & -- & -- & -- & --\\
\bottomrule
\end{tabular}
\end{table}

\paragraph{Vibration measures.}
For each HRNU of duration $T$ and axis $j$ we compute
\begin{gather}
a_{w,j}=\Bigl(\tfrac1T\textstyle\int_0^Ta_{w,j}^2\,dt\Bigr)^{1/2},\qquad
a_v=\bigl(a_{w,x}^2+a_{w,y}^2+a_{w,z}^2\bigr)^{1/2},\\
\mathrm{MTVV}_j=\max_{t_0}\Bigl(\tfrac1\tau\textstyle\int_{t_0-\tau}^{t_0}a_{w,j}^2\,dt\Bigr)^{1/2}\ (\tau=1~\mathrm{s}),\qquad
\mathrm{VDV}_j=\Bigl(\textstyle\int_0^Ta_{w,j}^4\,dt\Bigr)^{1/4},\qquad
\mathrm{CF}_j=\frac{\max_t|a_{w,j}(t)|}{a_{w,j}}.
\end{gather}
Following the standard's evaluation rule, $a_w$ is always reported.
MTVV and VDV are added, and the unit is flagged as shock-dominated, when
\begin{equation}
S_j=\ind\bigl[\mathrm{CF}_j>9\ \lor\ \mathrm{MTVV}_j/a_{w,j}>1.5\ \lor\ \mathrm{VDV}_j/(a_{w,j}T^{1/4})>1.75\bigr].
\end{equation}
For a zero-mean stationary signal $\mathrm{VDV}/(a_wT^{1/4})=\kappa^{1/4}$, where $\kappa$ is the kurtosis.
The VDV criterion is therefore the condition $\kappa>9.38$; Gaussian vibration gives a ratio of 1.32.
Comfort descriptors follow the informative comfort scale of the standard (Annex~C).
Its bands overlap, so a single label uses geometric-mean cut-points $\boldsymbol\theta=(0.315,0.561,0.894,1.414,2.236)~\mathrm{m\,s^{-2}}$.
ISO~2631-1 defines comfort at the seat--occupant interface, whereas a phone in a holder measures the cabin structure.
We therefore call these quantities \emph{weighted cabin acceleration} and validate them against a phone fixed at the seat rail (Section~\ref{sec:data}).
Below about 1~Hz, $a_{w,x}$ and $a_{w,y}$ mostly reflect braking and cornering, and we report them as driving rather than road measures.

\paragraph{Modal band ratios.}
From the Welch PSD of $a^v_z$ (Hann, 5.12~s segments, 50\% overlap) we compute energies in a body band $[0.8,2.5]$~Hz, a wheel-hop band centred on the vehicle's wheel-hop peak $\hat f_{wh}\in[6,20]$~Hz, and a high band $[16,40]$~Hz.
From these we form $R_{wb}=\log(E_{\mathrm{wheel}}/E_{\mathrm{body}})$ and $R_{hb}=\log(E_{\mathrm{high}}/E_{\mathrm{body}})$.
Through $\lambda=v/f$ these bands correspond to wavelengths from about 0.25~m to 12.5~m at 10~m\,s$^{-1}$, i.e.\ unevenness and megatexture, not pavement macrotexture (0.5--50~mm~\cite{iso13473}).
They describe how the vehicle's modes are excited, and we use them within a vehicle and for segmentation.

\paragraph{Impulsive events and maneuver context.}
We band-pass the vertical channel, $\tilde z=\mathrm{BP}_{[4,25]\,\mathrm{Hz}}(a^v_z)$, and apply the Teager--Kaiser energy operator~\cite{kaiser1990tkeo},
\begin{equation}
\Psi[n]=\tilde z[n]^2-\tilde z[n-1]\,\tilde z[n+1],\qquad E[n]=\bigl(h_{50\,\mathrm{ms}}*\max(\Psi,0)\bigr)[n].
\end{equation}
For $\tilde z=A\cos(\Omega n+\phi)$, $\Psi=A^2\sin^2\Omega$, so the operator responds to both amplitude and frequency.
Candidates are maximal runs where $E$ exceeds $\theta[n]=\operatorname{med}_W(E)+\kappa\cdot1.4826\operatorname{MAD}_W(E)$ over a 10~s centred window.
The adaptive threshold follows the background, which is much higher on cobblestone than on asphalt.
Runs closer than 0.3~s are merged, and runs longer than $D_{\max}=0.5~\mathrm{s}+3~\mathrm{m}/\bar v$ are passed to segmentation as rough stretches.
The maneuver state $m_e$ is one of \textsc{brake}, \textsc{accel}, \textsc{turn}, \textsc{swerve}, \textsc{stop-go} or \textsc{cruise}.
Longitudinal states come from GNSS ($a^G_{\mathrm{lon}}$) and lateral ones from the yaw rate ($v\,\omega^v_z$); \textsc{stop-go} requires $v<2~\mathrm{m\,s^{-1}}$ and a longitudinal jerk $|j_x|>J_{sg}$.
With $\rho_e=E_{[8,25]}(e)/E_{[0.5,4]}(e)$ the band-energy ratio of $a^v_z$ in the event window, each candidate is labelled
\begin{equation}
\ell(e)=
\begin{cases}
\text{\textsc{artifact}} & m_e\in\{\text{\textsc{stop-go}},\text{\textsc{accel}}\}\ \land\ \max_{n\in e}|j_x[n]|>J_x\ \land\ \rho_e<\rho_0,\\
\text{\textsc{defect}} & \lnot\,\text{\textsc{artifact}}\ \land\ \rho_e\ge\rho_0\ \land\ d_e\le D_{\max},\\
\text{\textsc{ambiguous}} & \text{otherwise}.
\end{cases}
\label{eq:gate}
\end{equation}
Braking and turning never veto a defect.
In our data they do not increase vertical vibration (Section~\ref{sec:l1eval}), and drivers slow down \emph{before} speed bumps, so a braking veto would remove true events.
Acceleration is gated because it is the only maneuver that adds a dashboard-only high-band component.
A swerve without a vertical candidate within $\pm2$~s yields a separate low-confidence \emph{avoidance} event.
A logistic model trained on held-out routes maps the event features to $p_{\mathrm{defect}}$.

% ---------------------------------------------------------------------
\subsection{Layer 2: Spatial Grounding and Segmentation}
\label{sec:layer2}

\paragraph{Map matching.}
We match GNSS fixes to OpenStreetMap with a hidden Markov model~\cite{newson2009hmm}.
The states are candidate projections $x_{k,i}$ within 50~m of fix $z_k$, the emissions are $\mathcal N(\lVert z_k-x_{k,i}\rVert_{gc};0,\sigma_z^2)$, and the transitions are $\beta^{-1}\exp(-|\lVert z_k-z_{k+1}\rVert_{gc}-\lVert x_{k,i}-x_{k+1,j}\rVert_{\mathrm{route}}|/\beta)$, decoded with Viterbi.
Forward--backward marginals give a per-fix match confidence.
The matched arc length $s(t)$ replaces $\int v\,dt$.
Each matched way contributes its \texttt{surface}, \texttt{smoothness}, \texttt{highway} and \texttt{name} tags, and nodes such as \texttt{traffic\_calming} and level crossings become candidate anchors.

\paragraph{Map-informed change-point detection.}
We resample Layer-1 features on a 5~m along-track grid, $\mathbf y_j=[\log\tilde a_{w,z},\log a_{w,z},R_{wb},R_{hb},\text{event rate},\log a_{w,y}]_j$, robust-standardized per trip.
Speed is excluded so that boundaries do not follow speed changes.
Segmentation solves
\begin{equation}
\min_{m,\,\tau_{1:m}}\ \sum_{i=1}^{m+1}C\bigl(\mathbf y_{\tau_{i-1}+1:\tau_i}\bigr)+\sum_{i=1}^{m}\beta(\tau_i),\qquad
C(\mathbf y_{a:b})=\sum_{j=a}^{b}\lVert\mathbf y_j-\bar{\mathbf y}_{a:b}\rVert^2_{\mathbf D^{-1}},
\label{eq:pelt}
\end{equation}
with a penalty that falls near OpenStreetMap transitions $b\in\mathcal B$ (surface-tag changes, road-class changes, junctions):
\begin{equation}
\beta(\tau)=\beta_0\bigl(1-\lambda\,\pi(\tau)\bigr),\qquad
\pi(\tau)=\max_{b\in\mathcal B}w_b\exp\!\Bigl(-\frac{(s_\tau-s_b)^2}{2\sigma_b^2}\Bigr),\qquad \lambda\in[0,1).
\end{equation}
Setting $\lambda=0$ recovers PELT~\cite{killick2012pelt}.
Let $F(0)=0$, $G(t)=F(t)+\beta(t)\ind[t>0]$ and $F(s)=\min_{t\in\mathcal R_s}\{G(t)+C(\mathbf y_{t+1:s})\}$.

\begin{lemma}[Pruning under location-dependent penalties]
\label{lem:prune}
Suppose $C(\mathbf y_{t+1:s})+C(\mathbf y_{s+1:T})\le C(\mathbf y_{t+1:T})$ for all $t<s<T$.
If $G(t)+C(\mathbf y_{t+1:s})\ge G(s)$, then $t$ is not the optimal last change point for any $T>s$ and can be removed from $\mathcal R$.
\end{lemma}
\begin{proof}
For $T>s$, $G(t)+C(\mathbf y_{t+1:T})\ge G(t)+C(\mathbf y_{t+1:s})+C(\mathbf y_{s+1:T})\ge G(s)+C(\mathbf y_{s+1:T})$.
\end{proof}

The usual condition $F(t)+C(\mathbf y_{t+1:s})\ge F(s)$ coincides with Lemma~\ref{lem:prune} only for a constant penalty.
With a lower penalty at map transitions it discards valid candidates.
In 300 random series it returned a sub-optimal partition every time, whereas the rule of Lemma~\ref{lem:prune} matched exhaustive optimal partitioning in all 300.
We choose $\beta_0=c\,d\ln n$ and $(c,\lambda)$ by leave-one-route-out validation, and report the penalty path~\cite{haynes2017crops}.

\paragraph{Arbitration and HRNUs.}
A change point within 25~m of a relevant map transition is snapped to it: the map fixes the location and the signal decides existence.
A map transition without a change point is kept only if a two-sample Hotelling test on $\mathbf y$ across it rejects at $\alpha=0.01$; otherwise every way split in the map would start a new unit.
Units shorter than 50~m are merged into their closer neighbour.
Each boundary is typed as map+signal, signal-only or map-only.
Each unit carries a map--signal disagreement score $1-P(\text{surface}=\text{OSM tag}\mid\mathbf y)$ when the map has a surface tag.
The resulting HRNUs, with their measures, events, anchors and confidences, are the input to Layer~3.

% ---------------------------------------------------------------------
\subsection{Layer 3: Kinetic--Spatial Tokenization and Grounded Generation}
\label{sec:layer3}

\paragraph{Claim ledger.}
Layer~3 first converts HRNUs and events into a ledger $\mathcal L$ of typed facts with identifiers, values, tolerances and confidences.
The types are severity, numeric, event, anchor, comparison, missing and metadata.
A \textsc{missing} fact records that a value is unavailable and why, e.g.\ a 48~s GNSS gap.
The ledger is the reference for both the run-time verifier and the evaluation.

\paragraph{Kinetic--Spatial Tokenizer.}
Algorithm~\ref{alg:kst} compiles the ledger into a capsule that fits a token budget $B$.
Each unit has a salience $\sigma_m$ and three serializations (omit, brief, full) with token costs $c_{m,\ell}$ measured with the target tokenizer.
Selection solves a multiple-choice knapsack exactly by dynamic programming:
\begin{equation}
V_m(b)=\max_{\ell\in\Lambda_m}\bigl\{V_{m-1}(b-c_{m,\ell})+\sigma_m\eta_\ell\bigr\},\qquad V_0(b)=0,
\label{eq:knap}
\end{equation}
where $\Lambda_m$ excludes \textsc{omit} for units in the must-mention set $\mathcal M$ and $\eta_{\textsc{omit}}=0$.
The cost is $O(3MB)$ for $M$ units.

\begin{algorithm}[t]
\caption{Kinetic--Spatial Tokenizer (KST)}
\label{alg:kst}
\begin{algorithmic}[1]
\Require HRNUs $\mathcal U=(u_1,\dots,u_M)$ in route order; events $\mathcal E$; anchors $\mathcal A$; optional history $\mathcal H$; audience $\pi$; budget $B$; tokenizer $\mathrm{tok}$
\Ensure capsule $\kappa$ with $|\mathrm{tok}(\kappa)|\le B$; ledger $\mathcal L$
\State $\mathcal L\gets\textsc{MetaFacts}(\text{channels},\text{map coverage},\pi)$
\For{$m=1,\dots,M$}
  \State $k_m\gets\textsc{IsoClass}(a_{v,m})$;\quad $\tilde k_m\gets\textsc{IsoClass}(\tilde a_{w,m})$ \Comment{felt vs.\ road}
  \State $\mathcal E_m\gets\{e\in\mathcal E: s(e)\in[s^0_m,s^1_m)\}$
  \State $\alpha_m\gets$ nearest anchor in $\mathcal A$ within $r_{\max}$ of $u_m$, or $\varnothing$
  \State $\sigma_m\gets w_1k_m+w_2\sum_{e\in\mathcal E_m}p_e\,\mathrm{sev}(e)+w_3\,\Delta_{\mathcal H}(u_m)+w_4\,\delta^{\mathrm{map}}_m$
  \State $\mathcal L\gets\mathcal L\cup\textsc{Facts}(u_m,k_m,\tilde k_m,\mathcal E_m,\alpha_m)$ \Comment{including \textsc{missing} facts}
  \For{$\ell\in\{\textsc{omit},\textsc{brief},\textsc{full}\}$}
    \State $c_{m,\ell}\gets|\mathrm{tok}(\textsc{Ser}_\ell(u_m))|$
  \EndFor
\EndFor
\State $\mathcal M\gets\{m:\exists e\in\mathcal E_m,\ p_e\,\mathrm{sev}(e)\ge\eta^\star\}$ \Comment{must-mention units}
\State $B'\gets B-|\mathrm{tok}(\textsc{Header})|-c_{\mathrm{rest}}$
\State $\mathbf x^\star\gets\textsc{McKnapsack}(\{c_{m,\ell}\},\{\sigma_m\eta_\ell\},B',\mathcal M)$ \Comment{Eq.~\eqref{eq:knap}}
\If{$\mathbf x^\star$ is infeasible}
  \State merge adjacent units outside $\mathcal M$ and repeat; if still infeasible, increase $B$ and log it
\EndIf
\State $\kappa\gets\textsc{Header}\ \|\ \textsc{Ser}_{x^\star_1}(u_1)\ \|\cdots\|\ \textsc{Ser}_{x^\star_M}(u_M)\ \|\ \textsc{Rest}(\{u_m:x^\star_m=\textsc{omit}\})$
\State \Return $\kappa,\ \mathcal L$
\end{algorithmic}
\end{algorithm}

\paragraph{Plan--ground--realize.}
Generation runs in four steps.
\emph{Plan:} the SLM emits a JSON plan listing which units and facts to mention, their order, a rhetorical relation for each (elaboration, contrast, cause, comparison with history) and a hedge.
Decoding is constrained by a grammar generated per capsule from the ledger identifiers~\cite{ggerganov_llamacpp}, so the plan is always well-formed and cannot reference a unit or event that does not exist.
\emph{Ground:} deterministic code expands each plan item into a fact frame with values, names and units taken from $\mathcal L$.
\emph{Realize:} the SLM writes text from the fact frames, either with placeholders for numbers and names that are filled afterwards, or as free text.
\emph{Verify:} claims extracted from the text (Section~\ref{sec:metrics}) are checked against $\mathcal L$.
A failing sentence is regenerated at most twice and then replaced by the template realization of the same plan item.
The SLM therefore controls selection, ordering, aggregation and phrasing, but not values.

\paragraph{Hedging.}
The plan assigns hedges from fact confidence $p$: assertive for $p\ge0.9$, \emph{likely} for $0.7\le p<0.9$, \emph{possibly} for $0.4\le p<0.7$; facts with $p<0.4$ are mentioned only when asked about.
The verifier rejects a sentence whose hedge is stronger than its confidence band allows.
 inside an acmart document.
% Preamble requirements (acmart already loads amsmath, amsthm, booktabs, xcolor):
%   \usepackage{algorithm}
%   \usepackage{algpseudocode}
% Bibliography keys used below (add to your .bib):
%   iso2631, iso8855, iso8608, iso13473, newson2009hmm, killick2012pelt,
%   haynes2017crops, kaiser1990tkeo, mohan2008nericell, eriksson2008pothole,
%   menegazzo2021pvs, ggerganov_llamacpp
% =====================================================================
\providecommand{\R}{\mathbb{R}}
\providecommand{\ind}{\mathbf{1}}
\providecommand{\TBD}[1][]{\textcolor{red}{[TBD#1]}}

\section{The RoadScribe System}
\label{sec:system}

RoadScribe turns a trip recorded by a commodity smartphone into a short narrative and answers follow-up questions about it.
The pipeline has three layers (Figure~\ref{fig:arch}).
Layer~1 converts inertial and GNSS samples into vehicle-frame ride measures and candidate events.
Layer~2 anchors these measures to the road network and partitions the trip into homogeneous road narrative units (HRNUs).
Layer~3 compiles the units into a token-budgeted capsule and a machine-checkable claim ledger, and generates text with an on-device small language model (SLM) under a plan--ground--realize protocol.

Four principles shape the design.
(P1)~Every number, place name and event in the output can be traced to a ledger fact.
(P2)~The language model decides what to say and how to say it, but never supplies a value; values come from the ledger.
(P3)~Uncertainty estimated in Layers~1--2 reaches the text as calibrated hedges.
(P4)~When verification fails, a deterministic realizer takes over, so the SLM can add coverage but cannot remove faithfulness.

\begin{figure}[t]
  \centering
  \fbox{\parbox{0.95\linewidth}{\centering\vspace{2em}Architecture figure placeholder: IMU+GNSS $\rightarrow$ Layer 1 (orientation, ISO weighting, events) $\rightarrow$ Layer 2 (map matching, MI-PELT, HRNUs) $\rightarrow$ Layer 3 (ledger, KST capsule, plan--ground--realize, verifier, template fallback) $\rightarrow$ narrative / answers\vspace{2em}}}
  \caption{RoadScribe architecture.}
  \label{fig:arch}
\end{figure}

% ---------------------------------------------------------------------
\subsection{Layer 1: Kinetic Signal Processing}
\label{sec:layer1}

\paragraph{Notation.}
The IMU delivers specific force $\mathbf f^b[n]\in\R^3$ and angular rate $\boldsymbol\omega^b[n]\in\R^3$ in the phone body frame $b$ at rate $f_s$ ($T_s=1/f_s$).
GNSS delivers speed $v_k$ at times $t_k$ (1~Hz).
The vehicle frame $v$ follows ISO~8855~\cite{iso8855} ($x$ forward, $y$ left, $z$ up).
We estimate $\mathbf R_{vb}\in SO(3)$ and the gravity-free acceleration
\begin{equation}
  \mathbf a^v[n]=\mathbf R_{vb}\,\mathbf f^b[n]-g\,\mathbf e_z .
  \label{eq:av}
\end{equation}
Speed is interpolated to IMU time with a monotone cubic (PCHIP), giving $v(t)$ and $a^G_{\mathrm{lon}}(t)=\dot v(t)$.

\paragraph{Mount episodes.}
A phone in a holder is static relative to the vehicle except when handled.
We mark candidate handling where the gyroscope reports rotation that the vehicle does not perform, $\lVert\bar{\boldsymbol\omega}^b[n]\rVert-|\dot\psi^G(t_n)|>\Omega_h$, with $\dot\psi^G$ the GNSS course rate.
A candidate becomes an episode boundary when the quasi-static gravity directions on its two sides differ by more than $\theta_h$.
Orientation is estimated per episode $e$.

\paragraph{Gravity and tilt.}
Let $\bar{\mathbf f}^b$ and $\bar{\boldsymbol\omega}^b$ be low-passed at 0.5~Hz. The quasi-static set of episode $e$ is
\begin{equation}
\mathcal Q_e=\bigl\{n:\ \bigl|\lVert\bar{\mathbf f}^b[n]\rVert-g\bigr|<\epsilon_g,\ \lVert\bar{\boldsymbol\omega}^b[n]\rVert<\epsilon_\omega,\ |a^G_{\mathrm{lon}}(t_n)|<\epsilon_a,\ v(t_n)>v_{\min}\bigr\},
\end{equation}
and the up direction is $\mathbf u_e=\hat{\mathbf g}_e/\lVert\hat{\mathbf g}_e\rVert$ with $\hat{\mathbf g}_e$ the geometric median of $\{\bar{\mathbf f}^b[n]:n\in\mathcal Q_e\}$, which is robust to outliers and equivariant to rotation.
Gating matters: an ungated low-pass estimate absorbs sustained acceleration and tilts by $\arctan(a/g)$, i.e.\ $17^\circ$ at $3~\mathrm{m\,s^{-2}}$.
With $\mathbf k=\mathbf u_e\times\mathbf e_z/\lVert\mathbf u_e\times\mathbf e_z\rVert$ and $\theta=\operatorname{atan2}(\lVert\mathbf u_e\times\mathbf e_z\rVert,\mathbf u_e^\top\mathbf e_z)$, Rodrigues' formula gives the leveling rotation
\begin{equation}
\mathbf R_1=\mathbf I+\sin\theta\,[\mathbf k]_\times+(1-\cos\theta)\,[\mathbf k]_\times^2 .
\end{equation}
The route-average road grade biases this tilt (a 5\% grade is $2.9^\circ$).
The resulting error is at most 0.13\% on the vertical channel, but up to $0.5~\mathrm{m\,s^{-2}}$ on the longitudinal channel.
We therefore take longitudinal maneuver thresholds from GNSS rather than from the IMU.

\paragraph{Heading.}
After leveling, the horizontal axes differ from the vehicle axes by an unknown yaw $\psi$.
The leveled yaw rate $\omega^\ell_z=(\mathbf R_1\bar{\boldsymbol\omega}^b)_z$ does not depend on $\psi$.
This yields two GNSS-referenced targets at every fix: the longitudinal acceleration $a^G_{\mathrm{lon},k}$ and the centripetal acceleration $a^G_{\mathrm{lat},k}=v_k\,\omega^\ell_z(t_k)$.
With the leveled horizontal specific force $\mathbf h_k\in\R^2$ and $\mathbf r_k=(a^G_{\mathrm{lon},k},a^G_{\mathrm{lat},k})^\top$,
\begin{equation}
\hat\psi=\arg\min_{\psi}\sum_k\bigl\lVert\mathbf R_2(\psi)\mathbf h_k-\mathbf r_k\bigr\rVert^2
=\operatorname{atan2}\bigl(M_{21}-M_{12},\ M_{11}+M_{22}\bigr),\qquad \mathbf M=\sum_k\mathbf h_k\mathbf r_k^\top,
\label{eq:yaw}
\end{equation}
where $\mathbf R_2(\psi)=\bigl[\begin{smallmatrix}\cos\psi&\sin\psi\\-\sin\psi&\cos\psi\end{smallmatrix}\bigr]$.
The closed form follows from maximizing $\operatorname{tr}(\mathbf R_2\mathbf M)$.
Only fixes with $v_k>3~\mathrm{m\,s^{-1}}$ are used.
Finally $\mathbf R_{vb}=\mathbf R_z(\hat\psi)\mathbf R_1$, and a bootstrap over fixes provides the orientation confidence passed to Layer~3.

\paragraph{Speed decoupling.}
Vibration grows with speed.
For a road whose elevation PSD follows ISO~8608~\cite{iso8608}, $G_d(n)=G_d(n_0)(n/n_0)^{-w}$, the temporal input PSD at speed $v$ scales as $v^{w-1}f^{-w}$.
For a fixed vehicle response the RMS response therefore scales as $v^{(w-1)/2}$.
Rather than assume $w$, we estimate the elasticity with the window-level model $\log a_{w,i}=\alpha_{j(i)}+\gamma\log v_i+\varepsilon_i$, where $j$ indexes homogeneous segments, and report
\begin{equation}
\tilde a_w=a_w\,(v_{\mathrm{ref}}/\bar v)^{\hat\gamma},\qquad \bar v\ge v_{\min},\ v_{\mathrm{ref}}=10~\mathrm{m\,s^{-1}}.
\end{equation}
The raw $a_w$ describes how the ride felt and $\tilde a_w$ describes the road; narratives keep the two apart.

\paragraph{Frequency weighting.}
ISO~2631-1~\cite{iso2631} weights seated whole-body vibration with $W_k$ (vertical) and $W_d$ (horizontal).
Each is a product of band-limiting filters, an acceleration--velocity transition and, for $W_k$, an upward step ($p=j2\pi f$, $\omega_i=2\pi f_i$, $Q_1=1/\sqrt2$):
\begin{align}
H_h(p)&=\frac{p^2}{p^2+\omega_1p/Q_1+\omega_1^2}, &
H_l(p)&=\frac{\omega_2^2}{p^2+\omega_2p/Q_1+\omega_2^2}, \nonumber\\
H_t(p)&=\frac{(\omega_4^2/\omega_3)(p+\omega_3)}{p^2+\omega_4p/Q_4+\omega_4^2}, &
H_s(p)&=\frac{p^2+\omega_5p/Q_5+\omega_5^2}{p^2+\omega_6p/Q_6+\omega_6^2},
\end{align}
with $W=H_hH_lH_tH_s$ and the parameters of Table~\ref{tab:iso}.
Narratives are produced after the trip, so we apply $W$ in the frequency domain, $a_w=\mathcal F^{-1}\{\mathcal F\{a\}\,W(j2\pi f)\}$ with 10~s of zero padding.
This reproduces the analog magnitude and phase exactly below $f_s/2$.
For streaming use, a bilinear transform per second-order section, pre-warped at each section's corner frequency, is an alternative.
At $f_s=100$~Hz its magnitude error is at most 0.72~dB ($W_k$) and 1.22~dB ($W_d$) up to 20~Hz, but it reaches 7--8~dB at 40~Hz and underestimates impulsive peaks by 3.8\%.
At $f_s=400$~Hz the error stays within 0.23~dB up to 40~Hz.
The 100~Hz band limit $H_l$ cannot be realized for $f_s\le200$~Hz; omitting it changes $W$ by less than 0.11~dB below 40~Hz.
At $f_s=100$~Hz our analysis band ends near 40~Hz, below the 80~Hz upper limit of the standard.

\begin{table}[t]
\caption{ISO 2631-1 frequency-weighting parameters (Annex A). Frequencies in Hz.}
\label{tab:iso}
\begin{tabular}{lccccccccc}
\toprule
 & $f_1$ & $f_2$ & $f_3$ & $f_4$ & $Q_4$ & $f_5$ & $Q_5$ & $f_6$ & $Q_6$\\
\midrule
$W_k$ & 0.4 & 100 & 12.5 & 12.5 & 0.63 & 2.37 & 0.91 & 3.35 & 0.91\\
$W_d$ & 0.4 & 100 & 2.0 & 2.0 & 0.63 & -- & -- & -- & --\\
\bottomrule
\end{tabular}
\end{table}

\paragraph{Vibration measures.}
For each HRNU of duration $T$ and axis $j$ we compute
\begin{gather}
a_{w,j}=\Bigl(\tfrac1T\textstyle\int_0^Ta_{w,j}^2\,dt\Bigr)^{1/2},\qquad
a_v=\bigl(a_{w,x}^2+a_{w,y}^2+a_{w,z}^2\bigr)^{1/2},\\
\mathrm{MTVV}_j=\max_{t_0}\Bigl(\tfrac1\tau\textstyle\int_{t_0-\tau}^{t_0}a_{w,j}^2\,dt\Bigr)^{1/2}\ (\tau=1~\mathrm{s}),\qquad
\mathrm{VDV}_j=\Bigl(\textstyle\int_0^Ta_{w,j}^4\,dt\Bigr)^{1/4},\qquad
\mathrm{CF}_j=\frac{\max_t|a_{w,j}(t)|}{a_{w,j}}.
\end{gather}
Following the standard's evaluation rule, $a_w$ is always reported.
MTVV and VDV are added, and the unit is flagged as shock-dominated, when
\begin{equation}
S_j=\ind\bigl[\mathrm{CF}_j>9\ \lor\ \mathrm{MTVV}_j/a_{w,j}>1.5\ \lor\ \mathrm{VDV}_j/(a_{w,j}T^{1/4})>1.75\bigr].
\end{equation}
For a zero-mean stationary signal $\mathrm{VDV}/(a_wT^{1/4})=\kappa^{1/4}$, where $\kappa$ is the kurtosis.
The VDV criterion is therefore the condition $\kappa>9.38$; Gaussian vibration gives a ratio of 1.32.
Comfort descriptors follow the informative comfort scale of the standard (Annex~C).
Its bands overlap, so a single label uses geometric-mean cut-points $\boldsymbol\theta=(0.315,0.561,0.894,1.414,2.236)~\mathrm{m\,s^{-2}}$.
ISO~2631-1 defines comfort at the seat--occupant interface, whereas a phone in a holder measures the cabin structure.
We therefore call these quantities \emph{weighted cabin acceleration} and validate them against a phone fixed at the seat rail (Section~\ref{sec:data}).
Below about 1~Hz, $a_{w,x}$ and $a_{w,y}$ mostly reflect braking and cornering, and we report them as driving rather than road measures.

\paragraph{Modal band ratios.}
From the Welch PSD of $a^v_z$ (Hann, 5.12~s segments, 50\% overlap) we compute energies in a body band $[0.8,2.5]$~Hz, a wheel-hop band centred on the vehicle's wheel-hop peak $\hat f_{wh}\in[6,20]$~Hz, and a high band $[16,40]$~Hz.
From these we form $R_{wb}=\log(E_{\mathrm{wheel}}/E_{\mathrm{body}})$ and $R_{hb}=\log(E_{\mathrm{high}}/E_{\mathrm{body}})$.
Through $\lambda=v/f$ these bands correspond to wavelengths from about 0.25~m to 12.5~m at 10~m\,s$^{-1}$, i.e.\ unevenness and megatexture, not pavement macrotexture (0.5--50~mm~\cite{iso13473}).
They describe how the vehicle's modes are excited, and we use them within a vehicle and for segmentation.

\paragraph{Impulsive events and maneuver context.}
We band-pass the vertical channel, $\tilde z=\mathrm{BP}_{[4,25]\,\mathrm{Hz}}(a^v_z)$, and apply the Teager--Kaiser energy operator~\cite{kaiser1990tkeo},
\begin{equation}
\Psi[n]=\tilde z[n]^2-\tilde z[n-1]\,\tilde z[n+1],\qquad E[n]=\bigl(h_{50\,\mathrm{ms}}*\max(\Psi,0)\bigr)[n].
\end{equation}
For $\tilde z=A\cos(\Omega n+\phi)$, $\Psi=A^2\sin^2\Omega$, so the operator responds to both amplitude and frequency.
Candidates are maximal runs where $E$ exceeds $\theta[n]=\operatorname{med}_W(E)+\kappa\cdot1.4826\operatorname{MAD}_W(E)$ over a 10~s centred window.
The adaptive threshold follows the background, which is much higher on cobblestone than on asphalt.
Runs closer than 0.3~s are merged, and runs longer than $D_{\max}=0.5~\mathrm{s}+3~\mathrm{m}/\bar v$ are passed to segmentation as rough stretches.
The maneuver state $m_e$ is one of \textsc{brake}, \textsc{accel}, \textsc{turn}, \textsc{swerve}, \textsc{stop-go} or \textsc{cruise}.
Longitudinal states come from GNSS ($a^G_{\mathrm{lon}}$) and lateral ones from the yaw rate ($v\,\omega^v_z$); \textsc{stop-go} requires $v<2~\mathrm{m\,s^{-1}}$ and a longitudinal jerk $|j_x|>J_{sg}$.
With $\rho_e=E_{[8,25]}(e)/E_{[0.5,4]}(e)$ the band-energy ratio of $a^v_z$ in the event window, each candidate is labelled
\begin{equation}
\ell(e)=
\begin{cases}
\text{\textsc{artifact}} & m_e\in\{\text{\textsc{stop-go}},\text{\textsc{accel}}\}\ \land\ \max_{n\in e}|j_x[n]|>J_x\ \land\ \rho_e<\rho_0,\\
\text{\textsc{defect}} & \lnot\,\text{\textsc{artifact}}\ \land\ \rho_e\ge\rho_0\ \land\ d_e\le D_{\max},\\
\text{\textsc{ambiguous}} & \text{otherwise}.
\end{cases}
\label{eq:gate}
\end{equation}
Braking and turning never veto a defect.
In our data they do not increase vertical vibration (Section~\ref{sec:l1eval}), and drivers slow down \emph{before} speed bumps, so a braking veto would remove true events.
Acceleration is gated because it is the only maneuver that adds a dashboard-only high-band component.
A swerve without a vertical candidate within $\pm2$~s yields a separate low-confidence \emph{avoidance} event.
A logistic model trained on held-out routes maps the event features to $p_{\mathrm{defect}}$.

% ---------------------------------------------------------------------
\subsection{Layer 2: Spatial Grounding and Segmentation}
\label{sec:layer2}

\paragraph{Map matching.}
We match GNSS fixes to OpenStreetMap with a hidden Markov model~\cite{newson2009hmm}.
The states are candidate projections $x_{k,i}$ within 50~m of fix $z_k$, the emissions are $\mathcal N(\lVert z_k-x_{k,i}\rVert_{gc};0,\sigma_z^2)$, and the transitions are $\beta^{-1}\exp(-|\lVert z_k-z_{k+1}\rVert_{gc}-\lVert x_{k,i}-x_{k+1,j}\rVert_{\mathrm{route}}|/\beta)$, decoded with Viterbi.
Forward--backward marginals give a per-fix match confidence.
The matched arc length $s(t)$ replaces $\int v\,dt$.
Each matched way contributes its \texttt{surface}, \texttt{smoothness}, \texttt{highway} and \texttt{name} tags, and nodes such as \texttt{traffic\_calming} and level crossings become candidate anchors.

\paragraph{Map-informed change-point detection.}
We resample Layer-1 features on a 5~m along-track grid, $\mathbf y_j=[\log\tilde a_{w,z},\log a_{w,z},R_{wb},R_{hb},\text{event rate},\log a_{w,y}]_j$, robust-standardized per trip.
Speed is excluded so that boundaries do not follow speed changes.
Segmentation solves
\begin{equation}
\min_{m,\,\tau_{1:m}}\ \sum_{i=1}^{m+1}C\bigl(\mathbf y_{\tau_{i-1}+1:\tau_i}\bigr)+\sum_{i=1}^{m}\beta(\tau_i),\qquad
C(\mathbf y_{a:b})=\sum_{j=a}^{b}\lVert\mathbf y_j-\bar{\mathbf y}_{a:b}\rVert^2_{\mathbf D^{-1}},
\label{eq:pelt}
\end{equation}
with a penalty that falls near OpenStreetMap transitions $b\in\mathcal B$ (surface-tag changes, road-class changes, junctions):
\begin{equation}
\beta(\tau)=\beta_0\bigl(1-\lambda\,\pi(\tau)\bigr),\qquad
\pi(\tau)=\max_{b\in\mathcal B}w_b\exp\!\Bigl(-\frac{(s_\tau-s_b)^2}{2\sigma_b^2}\Bigr),\qquad \lambda\in[0,1).
\end{equation}
Setting $\lambda=0$ recovers PELT~\cite{killick2012pelt}.
Let $F(0)=0$, $G(t)=F(t)+\beta(t)\ind[t>0]$ and $F(s)=\min_{t\in\mathcal R_s}\{G(t)+C(\mathbf y_{t+1:s})\}$.

\begin{lemma}[Pruning under location-dependent penalties]
\label{lem:prune}
Suppose $C(\mathbf y_{t+1:s})+C(\mathbf y_{s+1:T})\le C(\mathbf y_{t+1:T})$ for all $t<s<T$.
If $G(t)+C(\mathbf y_{t+1:s})\ge G(s)$, then $t$ is not the optimal last change point for any $T>s$ and can be removed from $\mathcal R$.
\end{lemma}
\begin{proof}
For $T>s$, $G(t)+C(\mathbf y_{t+1:T})\ge G(t)+C(\mathbf y_{t+1:s})+C(\mathbf y_{s+1:T})\ge G(s)+C(\mathbf y_{s+1:T})$.
\end{proof}

The usual condition $F(t)+C(\mathbf y_{t+1:s})\ge F(s)$ coincides with Lemma~\ref{lem:prune} only for a constant penalty.
With a lower penalty at map transitions it discards valid candidates.
In 300 random series it returned a sub-optimal partition every time, whereas the rule of Lemma~\ref{lem:prune} matched exhaustive optimal partitioning in all 300.
We choose $\beta_0=c\,d\ln n$ and $(c,\lambda)$ by leave-one-route-out validation, and report the penalty path~\cite{haynes2017crops}.

\paragraph{Arbitration and HRNUs.}
A change point within 25~m of a relevant map transition is snapped to it: the map fixes the location and the signal decides existence.
A map transition without a change point is kept only if a two-sample Hotelling test on $\mathbf y$ across it rejects at $\alpha=0.01$; otherwise every way split in the map would start a new unit.
Units shorter than 50~m are merged into their closer neighbour.
Each boundary is typed as map+signal, signal-only or map-only.
Each unit carries a map--signal disagreement score $1-P(\text{surface}=\text{OSM tag}\mid\mathbf y)$ when the map has a surface tag.
The resulting HRNUs, with their measures, events, anchors and confidences, are the input to Layer~3.

% ---------------------------------------------------------------------
\subsection{Layer 3: Kinetic--Spatial Tokenization and Grounded Generation}
\label{sec:layer3}

\paragraph{Claim ledger.}
Layer~3 first converts HRNUs and events into a ledger $\mathcal L$ of typed facts with identifiers, values, tolerances and confidences.
The types are severity, numeric, event, anchor, comparison, missing and metadata.
A \textsc{missing} fact records that a value is unavailable and why, e.g.\ a 48~s GNSS gap.
The ledger is the reference for both the run-time verifier and the evaluation.

\paragraph{Kinetic--Spatial Tokenizer.}
Algorithm~\ref{alg:kst} compiles the ledger into a capsule that fits a token budget $B$.
Each unit has a salience $\sigma_m$ and three serializations (omit, brief, full) with token costs $c_{m,\ell}$ measured with the target tokenizer.
Selection solves a multiple-choice knapsack exactly by dynamic programming:
\begin{equation}
V_m(b)=\max_{\ell\in\Lambda_m}\bigl\{V_{m-1}(b-c_{m,\ell})+\sigma_m\eta_\ell\bigr\},\qquad V_0(b)=0,
\label{eq:knap}
\end{equation}
where $\Lambda_m$ excludes \textsc{omit} for units in the must-mention set $\mathcal M$ and $\eta_{\textsc{omit}}=0$.
The cost is $O(3MB)$ for $M$ units.

\begin{algorithm}[t]
\caption{Kinetic--Spatial Tokenizer (KST)}
\label{alg:kst}
\begin{algorithmic}[1]
\Require HRNUs $\mathcal U=(u_1,\dots,u_M)$ in route order; events $\mathcal E$; anchors $\mathcal A$; optional history $\mathcal H$; audience $\pi$; budget $B$; tokenizer $\mathrm{tok}$
\Ensure capsule $\kappa$ with $|\mathrm{tok}(\kappa)|\le B$; ledger $\mathcal L$
\State $\mathcal L\gets\textsc{MetaFacts}(\text{channels},\text{map coverage},\pi)$
\For{$m=1,\dots,M$}
  \State $k_m\gets\textsc{IsoClass}(a_{v,m})$;\quad $\tilde k_m\gets\textsc{IsoClass}(\tilde a_{w,m})$ \Comment{felt vs.\ road}
  \State $\mathcal E_m\gets\{e\in\mathcal E: s(e)\in[s^0_m,s^1_m)\}$
  \State $\alpha_m\gets$ nearest anchor in $\mathcal A$ within $r_{\max}$ of $u_m$, or $\varnothing$
  \State $\sigma_m\gets w_1k_m+w_2\sum_{e\in\mathcal E_m}p_e\,\mathrm{sev}(e)+w_3\,\Delta_{\mathcal H}(u_m)+w_4\,\delta^{\mathrm{map}}_m$
  \State $\mathcal L\gets\mathcal L\cup\textsc{Facts}(u_m,k_m,\tilde k_m,\mathcal E_m,\alpha_m)$ \Comment{including \textsc{missing} facts}
  \For{$\ell\in\{\textsc{omit},\textsc{brief},\textsc{full}\}$}
    \State $c_{m,\ell}\gets|\mathrm{tok}(\textsc{Ser}_\ell(u_m))|$
  \EndFor
\EndFor
\State $\mathcal M\gets\{m:\exists e\in\mathcal E_m,\ p_e\,\mathrm{sev}(e)\ge\eta^\star\}$ \Comment{must-mention units}
\State $B'\gets B-|\mathrm{tok}(\textsc{Header})|-c_{\mathrm{rest}}$
\State $\mathbf x^\star\gets\textsc{McKnapsack}(\{c_{m,\ell}\},\{\sigma_m\eta_\ell\},B',\mathcal M)$ \Comment{Eq.~\eqref{eq:knap}}
\If{$\mathbf x^\star$ is infeasible}
  \State merge adjacent units outside $\mathcal M$ and repeat; if still infeasible, increase $B$ and log it
\EndIf
\State $\kappa\gets\textsc{Header}\ \|\ \textsc{Ser}_{x^\star_1}(u_1)\ \|\cdots\|\ \textsc{Ser}_{x^\star_M}(u_M)\ \|\ \textsc{Rest}(\{u_m:x^\star_m=\textsc{omit}\})$
\State \Return $\kappa,\ \mathcal L$
\end{algorithmic}
\end{algorithm}

\paragraph{Plan--ground--realize.}
Generation runs in four steps.
\emph{Plan:} the SLM emits a JSON plan listing which units and facts to mention, their order, a rhetorical relation for each (elaboration, contrast, cause, comparison with history) and a hedge.
Decoding is constrained by a grammar generated per capsule from the ledger identifiers~\cite{ggerganov_llamacpp}, so the plan is always well-formed and cannot reference a unit or event that does not exist.
\emph{Ground:} deterministic code expands each plan item into a fact frame with values, names and units taken from $\mathcal L$.
\emph{Realize:} the SLM writes text from the fact frames, either with placeholders for numbers and names that are filled afterwards, or as free text.
\emph{Verify:} claims extracted from the text (Section~\ref{sec:metrics}) are checked against $\mathcal L$.
A failing sentence is regenerated at most twice and then replaced by the template realization of the same plan item.
The SLM therefore controls selection, ordering, aggregation and phrasing, but not values.

\paragraph{Hedging.}
The plan assigns hedges from fact confidence $p$: assertive for $p\ge0.9$, \emph{likely} for $0.7\le p<0.9$, \emph{possibly} for $0.4\le p<0.7$; facts with $p<0.4$ are mentioned only when asked about.
The verifier rejects a sentence whose hedge is stronger than its confidence band allows.
 inside an acmart document.
% Preamble requirements (acmart already loads amsmath, amsthm, booktabs, xcolor):
%   \usepackage{algorithm}
%   \usepackage{algpseudocode}
% Bibliography keys used below (add to your .bib):
%   iso2631, iso8855, iso8608, iso13473, newson2009hmm, killick2012pelt,
%   haynes2017crops, kaiser1990tkeo, mohan2008nericell, eriksson2008pothole,
%   menegazzo2021pvs, ggerganov_llamacpp
% =====================================================================
\providecommand{\R}{\mathbb{R}}
\providecommand{\ind}{\mathbf{1}}
\providecommand{\TBD}[1][]{\textcolor{red}{[TBD#1]}}

\section{The RoadScribe System}
\label{sec:system}

RoadScribe turns a trip recorded by a commodity smartphone into a short narrative and answers follow-up questions about it.
The pipeline has three layers (Figure~\ref{fig:arch}).
Layer~1 converts inertial and GNSS samples into vehicle-frame ride measures and candidate events.
Layer~2 anchors these measures to the road network and partitions the trip into homogeneous road narrative units (HRNUs).
Layer~3 compiles the units into a token-budgeted capsule and a machine-checkable claim ledger, and generates text with an on-device small language model (SLM) under a plan--ground--realize protocol.

Four principles shape the design.
(P1)~Every number, place name and event in the output can be traced to a ledger fact.
(P2)~The language model decides what to say and how to say it, but never supplies a value; values come from the ledger.
(P3)~Uncertainty estimated in Layers~1--2 reaches the text as calibrated hedges.
(P4)~When verification fails, a deterministic realizer takes over, so the SLM can add coverage but cannot remove faithfulness.

\begin{figure}[t]
  \centering
  \fbox{\parbox{0.95\linewidth}{\centering\vspace{2em}Architecture figure placeholder: IMU+GNSS $\rightarrow$ Layer 1 (orientation, ISO weighting, events) $\rightarrow$ Layer 2 (map matching, MI-PELT, HRNUs) $\rightarrow$ Layer 3 (ledger, KST capsule, plan--ground--realize, verifier, template fallback) $\rightarrow$ narrative / answers\vspace{2em}}}
  \caption{RoadScribe architecture.}
  \label{fig:arch}
\end{figure}

% ---------------------------------------------------------------------
\subsection{Layer 1: Kinetic Signal Processing}
\label{sec:layer1}

\paragraph{Notation.}
The IMU delivers specific force $\mathbf f^b[n]\in\R^3$ and angular rate $\boldsymbol\omega^b[n]\in\R^3$ in the phone body frame $b$ at rate $f_s$ ($T_s=1/f_s$).
GNSS delivers speed $v_k$ at times $t_k$ (1~Hz).
The vehicle frame $v$ follows ISO~8855~\cite{iso8855} ($x$ forward, $y$ left, $z$ up).
We estimate $\mathbf R_{vb}\in SO(3)$ and the gravity-free acceleration
\begin{equation}
  \mathbf a^v[n]=\mathbf R_{vb}\,\mathbf f^b[n]-g\,\mathbf e_z .
  \label{eq:av}
\end{equation}
Speed is interpolated to IMU time with a monotone cubic (PCHIP), giving $v(t)$ and $a^G_{\mathrm{lon}}(t)=\dot v(t)$.

\paragraph{Mount episodes.}
A phone in a holder is static relative to the vehicle except when handled.
We mark candidate handling where the gyroscope reports rotation that the vehicle does not perform, $\lVert\bar{\boldsymbol\omega}^b[n]\rVert-|\dot\psi^G(t_n)|>\Omega_h$, with $\dot\psi^G$ the GNSS course rate.
A candidate becomes an episode boundary when the quasi-static gravity directions on its two sides differ by more than $\theta_h$.
Orientation is estimated per episode $e$.

\paragraph{Gravity and tilt.}
Let $\bar{\mathbf f}^b$ and $\bar{\boldsymbol\omega}^b$ be low-passed at 0.5~Hz. The quasi-static set of episode $e$ is
\begin{equation}
\mathcal Q_e=\bigl\{n:\ \bigl|\lVert\bar{\mathbf f}^b[n]\rVert-g\bigr|<\epsilon_g,\ \lVert\bar{\boldsymbol\omega}^b[n]\rVert<\epsilon_\omega,\ |a^G_{\mathrm{lon}}(t_n)|<\epsilon_a,\ v(t_n)>v_{\min}\bigr\},
\end{equation}
and the up direction is $\mathbf u_e=\hat{\mathbf g}_e/\lVert\hat{\mathbf g}_e\rVert$ with $\hat{\mathbf g}_e$ the geometric median of $\{\bar{\mathbf f}^b[n]:n\in\mathcal Q_e\}$, which is robust to outliers and equivariant to rotation.
Gating matters: an ungated low-pass estimate absorbs sustained acceleration and tilts by $\arctan(a/g)$, i.e.\ $17^\circ$ at $3~\mathrm{m\,s^{-2}}$.
With $\mathbf k=\mathbf u_e\times\mathbf e_z/\lVert\mathbf u_e\times\mathbf e_z\rVert$ and $\theta=\operatorname{atan2}(\lVert\mathbf u_e\times\mathbf e_z\rVert,\mathbf u_e^\top\mathbf e_z)$, Rodrigues' formula gives the leveling rotation
\begin{equation}
\mathbf R_1=\mathbf I+\sin\theta\,[\mathbf k]_\times+(1-\cos\theta)\,[\mathbf k]_\times^2 .
\end{equation}
The route-average road grade biases this tilt (a 5\% grade is $2.9^\circ$).
The resulting error is at most 0.13\% on the vertical channel, but up to $0.5~\mathrm{m\,s^{-2}}$ on the longitudinal channel.
We therefore take longitudinal maneuver thresholds from GNSS rather than from the IMU.

\paragraph{Heading.}
After leveling, the horizontal axes differ from the vehicle axes by an unknown yaw $\psi$.
The leveled yaw rate $\omega^\ell_z=(\mathbf R_1\bar{\boldsymbol\omega}^b)_z$ does not depend on $\psi$.
This yields two GNSS-referenced targets at every fix: the longitudinal acceleration $a^G_{\mathrm{lon},k}$ and the centripetal acceleration $a^G_{\mathrm{lat},k}=v_k\,\omega^\ell_z(t_k)$.
With the leveled horizontal specific force $\mathbf h_k\in\R^2$ and $\mathbf r_k=(a^G_{\mathrm{lon},k},a^G_{\mathrm{lat},k})^\top$,
\begin{equation}
\hat\psi=\arg\min_{\psi}\sum_k\bigl\lVert\mathbf R_2(\psi)\mathbf h_k-\mathbf r_k\bigr\rVert^2
=\operatorname{atan2}\bigl(M_{21}-M_{12},\ M_{11}+M_{22}\bigr),\qquad \mathbf M=\sum_k\mathbf h_k\mathbf r_k^\top,
\label{eq:yaw}
\end{equation}
where $\mathbf R_2(\psi)=\bigl[\begin{smallmatrix}\cos\psi&\sin\psi\\-\sin\psi&\cos\psi\end{smallmatrix}\bigr]$.
The closed form follows from maximizing $\operatorname{tr}(\mathbf R_2\mathbf M)$.
Only fixes with $v_k>3~\mathrm{m\,s^{-1}}$ are used.
Finally $\mathbf R_{vb}=\mathbf R_z(\hat\psi)\mathbf R_1$, and a bootstrap over fixes provides the orientation confidence passed to Layer~3.

\paragraph{Speed decoupling.}
Vibration grows with speed.
For a road whose elevation PSD follows ISO~8608~\cite{iso8608}, $G_d(n)=G_d(n_0)(n/n_0)^{-w}$, the temporal input PSD at speed $v$ scales as $v^{w-1}f^{-w}$.
For a fixed vehicle response the RMS response therefore scales as $v^{(w-1)/2}$.
Rather than assume $w$, we estimate the elasticity with the window-level model $\log a_{w,i}=\alpha_{j(i)}+\gamma\log v_i+\varepsilon_i$, where $j$ indexes homogeneous segments, and report
\begin{equation}
\tilde a_w=a_w\,(v_{\mathrm{ref}}/\bar v)^{\hat\gamma},\qquad \bar v\ge v_{\min},\ v_{\mathrm{ref}}=10~\mathrm{m\,s^{-1}}.
\end{equation}
The raw $a_w$ describes how the ride felt and $\tilde a_w$ describes the road; narratives keep the two apart.

\paragraph{Frequency weighting.}
ISO~2631-1~\cite{iso2631} weights seated whole-body vibration with $W_k$ (vertical) and $W_d$ (horizontal).
Each is a product of band-limiting filters, an acceleration--velocity transition and, for $W_k$, an upward step ($p=j2\pi f$, $\omega_i=2\pi f_i$, $Q_1=1/\sqrt2$):
\begin{align}
H_h(p)&=\frac{p^2}{p^2+\omega_1p/Q_1+\omega_1^2}, &
H_l(p)&=\frac{\omega_2^2}{p^2+\omega_2p/Q_1+\omega_2^2}, \nonumber\\
H_t(p)&=\frac{(\omega_4^2/\omega_3)(p+\omega_3)}{p^2+\omega_4p/Q_4+\omega_4^2}, &
H_s(p)&=\frac{p^2+\omega_5p/Q_5+\omega_5^2}{p^2+\omega_6p/Q_6+\omega_6^2},
\end{align}
with $W=H_hH_lH_tH_s$ and the parameters of Table~\ref{tab:iso}.
Narratives are produced after the trip, so we apply $W$ in the frequency domain, $a_w=\mathcal F^{-1}\{\mathcal F\{a\}\,W(j2\pi f)\}$ with 10~s of zero padding.
This reproduces the analog magnitude and phase exactly below $f_s/2$.
For streaming use, a bilinear transform per second-order section, pre-warped at each section's corner frequency, is an alternative.
At $f_s=100$~Hz its magnitude error is at most 0.72~dB ($W_k$) and 1.22~dB ($W_d$) up to 20~Hz, but it reaches 7--8~dB at 40~Hz and underestimates impulsive peaks by 3.8\%.
At $f_s=400$~Hz the error stays within 0.23~dB up to 40~Hz.
The 100~Hz band limit $H_l$ cannot be realized for $f_s\le200$~Hz; omitting it changes $W$ by less than 0.11~dB below 40~Hz.
At $f_s=100$~Hz our analysis band ends near 40~Hz, below the 80~Hz upper limit of the standard.

\begin{table}[t]
\caption{ISO 2631-1 frequency-weighting parameters (Annex A). Frequencies in Hz.}
\label{tab:iso}
\begin{tabular}{lccccccccc}
\toprule
 & $f_1$ & $f_2$ & $f_3$ & $f_4$ & $Q_4$ & $f_5$ & $Q_5$ & $f_6$ & $Q_6$\\
\midrule
$W_k$ & 0.4 & 100 & 12.5 & 12.5 & 0.63 & 2.37 & 0.91 & 3.35 & 0.91\\
$W_d$ & 0.4 & 100 & 2.0 & 2.0 & 0.63 & -- & -- & -- & --\\
\bottomrule
\end{tabular}
\end{table}

\paragraph{Vibration measures.}
For each HRNU of duration $T$ and axis $j$ we compute
\begin{gather}
a_{w,j}=\Bigl(\tfrac1T\textstyle\int_0^Ta_{w,j}^2\,dt\Bigr)^{1/2},\qquad
a_v=\bigl(a_{w,x}^2+a_{w,y}^2+a_{w,z}^2\bigr)^{1/2},\\
\mathrm{MTVV}_j=\max_{t_0}\Bigl(\tfrac1\tau\textstyle\int_{t_0-\tau}^{t_0}a_{w,j}^2\,dt\Bigr)^{1/2}\ (\tau=1~\mathrm{s}),\qquad
\mathrm{VDV}_j=\Bigl(\textstyle\int_0^Ta_{w,j}^4\,dt\Bigr)^{1/4},\qquad
\mathrm{CF}_j=\frac{\max_t|a_{w,j}(t)|}{a_{w,j}}.
\end{gather}
Following the standard's evaluation rule, $a_w$ is always reported.
MTVV and VDV are added, and the unit is flagged as shock-dominated, when
\begin{equation}
S_j=\ind\bigl[\mathrm{CF}_j>9\ \lor\ \mathrm{MTVV}_j/a_{w,j}>1.5\ \lor\ \mathrm{VDV}_j/(a_{w,j}T^{1/4})>1.75\bigr].
\end{equation}
For a zero-mean stationary signal $\mathrm{VDV}/(a_wT^{1/4})=\kappa^{1/4}$, where $\kappa$ is the kurtosis.
The VDV criterion is therefore the condition $\kappa>9.38$; Gaussian vibration gives a ratio of 1.32.
Comfort descriptors follow the informative comfort scale of the standard (Annex~C).
Its bands overlap, so a single label uses geometric-mean cut-points $\boldsymbol\theta=(0.315,0.561,0.894,1.414,2.236)~\mathrm{m\,s^{-2}}$.
ISO~2631-1 defines comfort at the seat--occupant interface, whereas a phone in a holder measures the cabin structure.
We therefore call these quantities \emph{weighted cabin acceleration} and validate them against a phone fixed at the seat rail (Section~\ref{sec:data}).
Below about 1~Hz, $a_{w,x}$ and $a_{w,y}$ mostly reflect braking and cornering, and we report them as driving rather than road measures.

\paragraph{Modal band ratios.}
From the Welch PSD of $a^v_z$ (Hann, 5.12~s segments, 50\% overlap) we compute energies in a body band $[0.8,2.5]$~Hz, a wheel-hop band centred on the vehicle's wheel-hop peak $\hat f_{wh}\in[6,20]$~Hz, and a high band $[16,40]$~Hz.
From these we form $R_{wb}=\log(E_{\mathrm{wheel}}/E_{\mathrm{body}})$ and $R_{hb}=\log(E_{\mathrm{high}}/E_{\mathrm{body}})$.
Through $\lambda=v/f$ these bands correspond to wavelengths from about 0.25~m to 12.5~m at 10~m\,s$^{-1}$, i.e.\ unevenness and megatexture, not pavement macrotexture (0.5--50~mm~\cite{iso13473}).
They describe how the vehicle's modes are excited, and we use them within a vehicle and for segmentation.

\paragraph{Impulsive events and maneuver context.}
We band-pass the vertical channel, $\tilde z=\mathrm{BP}_{[4,25]\,\mathrm{Hz}}(a^v_z)$, and apply the Teager--Kaiser energy operator~\cite{kaiser1990tkeo},
\begin{equation}
\Psi[n]=\tilde z[n]^2-\tilde z[n-1]\,\tilde z[n+1],\qquad E[n]=\bigl(h_{50\,\mathrm{ms}}*\max(\Psi,0)\bigr)[n].
\end{equation}
For $\tilde z=A\cos(\Omega n+\phi)$, $\Psi=A^2\sin^2\Omega$, so the operator responds to both amplitude and frequency.
Candidates are maximal runs where $E$ exceeds $\theta[n]=\operatorname{med}_W(E)+\kappa\cdot1.4826\operatorname{MAD}_W(E)$ over a 10~s centred window.
The adaptive threshold follows the background, which is much higher on cobblestone than on asphalt.
Runs closer than 0.3~s are merged, and runs longer than $D_{\max}=0.5~\mathrm{s}+3~\mathrm{m}/\bar v$ are passed to segmentation as rough stretches.
The maneuver state $m_e$ is one of \textsc{brake}, \textsc{accel}, \textsc{turn}, \textsc{swerve}, \textsc{stop-go} or \textsc{cruise}.
Longitudinal states come from GNSS ($a^G_{\mathrm{lon}}$) and lateral ones from the yaw rate ($v\,\omega^v_z$); \textsc{stop-go} requires $v<2~\mathrm{m\,s^{-1}}$ and a longitudinal jerk $|j_x|>J_{sg}$.
With $\rho_e=E_{[8,25]}(e)/E_{[0.5,4]}(e)$ the band-energy ratio of $a^v_z$ in the event window, each candidate is labelled
\begin{equation}
\ell(e)=
\begin{cases}
\text{\textsc{artifact}} & m_e\in\{\text{\textsc{stop-go}},\text{\textsc{accel}}\}\ \land\ \max_{n\in e}|j_x[n]|>J_x\ \land\ \rho_e<\rho_0,\\
\text{\textsc{defect}} & \lnot\,\text{\textsc{artifact}}\ \land\ \rho_e\ge\rho_0\ \land\ d_e\le D_{\max},\\
\text{\textsc{ambiguous}} & \text{otherwise}.
\end{cases}
\label{eq:gate}
\end{equation}
Braking and turning never veto a defect.
In our data they do not increase vertical vibration (Section~\ref{sec:l1eval}), and drivers slow down \emph{before} speed bumps, so a braking veto would remove true events.
Acceleration is gated because it is the only maneuver that adds a dashboard-only high-band component.
A swerve without a vertical candidate within $\pm2$~s yields a separate low-confidence \emph{avoidance} event.
A logistic model trained on held-out routes maps the event features to $p_{\mathrm{defect}}$.

% ---------------------------------------------------------------------
\subsection{Layer 2: Spatial Grounding and Segmentation}
\label{sec:layer2}

\paragraph{Map matching.}
We match GNSS fixes to OpenStreetMap with a hidden Markov model~\cite{newson2009hmm}.
The states are candidate projections $x_{k,i}$ within 50~m of fix $z_k$, the emissions are $\mathcal N(\lVert z_k-x_{k,i}\rVert_{gc};0,\sigma_z^2)$, and the transitions are $\beta^{-1}\exp(-|\lVert z_k-z_{k+1}\rVert_{gc}-\lVert x_{k,i}-x_{k+1,j}\rVert_{\mathrm{route}}|/\beta)$, decoded with Viterbi.
Forward--backward marginals give a per-fix match confidence.
The matched arc length $s(t)$ replaces $\int v\,dt$.
Each matched way contributes its \texttt{surface}, \texttt{smoothness}, \texttt{highway} and \texttt{name} tags, and nodes such as \texttt{traffic\_calming} and level crossings become candidate anchors.

\paragraph{Map-informed change-point detection.}
We resample Layer-1 features on a 5~m along-track grid, $\mathbf y_j=[\log\tilde a_{w,z},\log a_{w,z},R_{wb},R_{hb},\text{event rate},\log a_{w,y}]_j$, robust-standardized per trip.
Speed is excluded so that boundaries do not follow speed changes.
Segmentation solves
\begin{equation}
\min_{m,\,\tau_{1:m}}\ \sum_{i=1}^{m+1}C\bigl(\mathbf y_{\tau_{i-1}+1:\tau_i}\bigr)+\sum_{i=1}^{m}\beta(\tau_i),\qquad
C(\mathbf y_{a:b})=\sum_{j=a}^{b}\lVert\mathbf y_j-\bar{\mathbf y}_{a:b}\rVert^2_{\mathbf D^{-1}},
\label{eq:pelt}
\end{equation}
with a penalty that falls near OpenStreetMap transitions $b\in\mathcal B$ (surface-tag changes, road-class changes, junctions):
\begin{equation}
\beta(\tau)=\beta_0\bigl(1-\lambda\,\pi(\tau)\bigr),\qquad
\pi(\tau)=\max_{b\in\mathcal B}w_b\exp\!\Bigl(-\frac{(s_\tau-s_b)^2}{2\sigma_b^2}\Bigr),\qquad \lambda\in[0,1).
\end{equation}
Setting $\lambda=0$ recovers PELT~\cite{killick2012pelt}.
Let $F(0)=0$, $G(t)=F(t)+\beta(t)\ind[t>0]$ and $F(s)=\min_{t\in\mathcal R_s}\{G(t)+C(\mathbf y_{t+1:s})\}$.

\begin{lemma}[Pruning under location-dependent penalties]
\label{lem:prune}
Suppose $C(\mathbf y_{t+1:s})+C(\mathbf y_{s+1:T})\le C(\mathbf y_{t+1:T})$ for all $t<s<T$.
If $G(t)+C(\mathbf y_{t+1:s})\ge G(s)$, then $t$ is not the optimal last change point for any $T>s$ and can be removed from $\mathcal R$.
\end{lemma}
\begin{proof}
For $T>s$, $G(t)+C(\mathbf y_{t+1:T})\ge G(t)+C(\mathbf y_{t+1:s})+C(\mathbf y_{s+1:T})\ge G(s)+C(\mathbf y_{s+1:T})$.
\end{proof}

The usual condition $F(t)+C(\mathbf y_{t+1:s})\ge F(s)$ coincides with Lemma~\ref{lem:prune} only for a constant penalty.
With a lower penalty at map transitions it discards valid candidates.
In 300 random series it returned a sub-optimal partition every time, whereas the rule of Lemma~\ref{lem:prune} matched exhaustive optimal partitioning in all 300.
We choose $\beta_0=c\,d\ln n$ and $(c,\lambda)$ by leave-one-route-out validation, and report the penalty path~\cite{haynes2017crops}.

\paragraph{Arbitration and HRNUs.}
A change point within 25~m of a relevant map transition is snapped to it: the map fixes the location and the signal decides existence.
A map transition without a change point is kept only if a two-sample Hotelling test on $\mathbf y$ across it rejects at $\alpha=0.01$; otherwise every way split in the map would start a new unit.
Units shorter than 50~m are merged into their closer neighbour.
Each boundary is typed as map+signal, signal-only or map-only.
Each unit carries a map--signal disagreement score $1-P(\text{surface}=\text{OSM tag}\mid\mathbf y)$ when the map has a surface tag.
The resulting HRNUs, with their measures, events, anchors and confidences, are the input to Layer~3.

% ---------------------------------------------------------------------
\subsection{Layer 3: Kinetic--Spatial Tokenization and Grounded Generation}
\label{sec:layer3}

\paragraph{Claim ledger.}
Layer~3 first converts HRNUs and events into a ledger $\mathcal L$ of typed facts with identifiers, values, tolerances and confidences.
The types are severity, numeric, event, anchor, comparison, missing and metadata.
A \textsc{missing} fact records that a value is unavailable and why, e.g.\ a 48~s GNSS gap.
The ledger is the reference for both the run-time verifier and the evaluation.

\paragraph{Kinetic--Spatial Tokenizer.}
Algorithm~\ref{alg:kst} compiles the ledger into a capsule that fits a token budget $B$.
Each unit has a salience $\sigma_m$ and three serializations (omit, brief, full) with token costs $c_{m,\ell}$ measured with the target tokenizer.
Selection solves a multiple-choice knapsack exactly by dynamic programming:
\begin{equation}
V_m(b)=\max_{\ell\in\Lambda_m}\bigl\{V_{m-1}(b-c_{m,\ell})+\sigma_m\eta_\ell\bigr\},\qquad V_0(b)=0,
\label{eq:knap}
\end{equation}
where $\Lambda_m$ excludes \textsc{omit} for units in the must-mention set $\mathcal M$ and $\eta_{\textsc{omit}}=0$.
The cost is $O(3MB)$ for $M$ units.

\begin{algorithm}[t]
\caption{Kinetic--Spatial Tokenizer (KST)}
\label{alg:kst}
\begin{algorithmic}[1]
\Require HRNUs $\mathcal U=(u_1,\dots,u_M)$ in route order; events $\mathcal E$; anchors $\mathcal A$; optional history $\mathcal H$; audience $\pi$; budget $B$; tokenizer $\mathrm{tok}$
\Ensure capsule $\kappa$ with $|\mathrm{tok}(\kappa)|\le B$; ledger $\mathcal L$
\State $\mathcal L\gets\textsc{MetaFacts}(\text{channels},\text{map coverage},\pi)$
\For{$m=1,\dots,M$}
  \State $k_m\gets\textsc{IsoClass}(a_{v,m})$;\quad $\tilde k_m\gets\textsc{IsoClass}(\tilde a_{w,m})$ \Comment{felt vs.\ road}
  \State $\mathcal E_m\gets\{e\in\mathcal E: s(e)\in[s^0_m,s^1_m)\}$
  \State $\alpha_m\gets$ nearest anchor in $\mathcal A$ within $r_{\max}$ of $u_m$, or $\varnothing$
  \State $\sigma_m\gets w_1k_m+w_2\sum_{e\in\mathcal E_m}p_e\,\mathrm{sev}(e)+w_3\,\Delta_{\mathcal H}(u_m)+w_4\,\delta^{\mathrm{map}}_m$
  \State $\mathcal L\gets\mathcal L\cup\textsc{Facts}(u_m,k_m,\tilde k_m,\mathcal E_m,\alpha_m)$ \Comment{including \textsc{missing} facts}
  \For{$\ell\in\{\textsc{omit},\textsc{brief},\textsc{full}\}$}
    \State $c_{m,\ell}\gets|\mathrm{tok}(\textsc{Ser}_\ell(u_m))|$
  \EndFor
\EndFor
\State $\mathcal M\gets\{m:\exists e\in\mathcal E_m,\ p_e\,\mathrm{sev}(e)\ge\eta^\star\}$ \Comment{must-mention units}
\State $B'\gets B-|\mathrm{tok}(\textsc{Header})|-c_{\mathrm{rest}}$
\State $\mathbf x^\star\gets\textsc{McKnapsack}(\{c_{m,\ell}\},\{\sigma_m\eta_\ell\},B',\mathcal M)$ \Comment{Eq.~\eqref{eq:knap}}
\If{$\mathbf x^\star$ is infeasible}
  \State merge adjacent units outside $\mathcal M$ and repeat; if still infeasible, increase $B$ and log it
\EndIf
\State $\kappa\gets\textsc{Header}\ \|\ \textsc{Ser}_{x^\star_1}(u_1)\ \|\cdots\|\ \textsc{Ser}_{x^\star_M}(u_M)\ \|\ \textsc{Rest}(\{u_m:x^\star_m=\textsc{omit}\})$
\State \Return $\kappa,\ \mathcal L$
\end{algorithmic}
\end{algorithm}

\paragraph{Plan--ground--realize.}
Generation runs in four steps.
\emph{Plan:} the SLM emits a JSON plan listing which units and facts to mention, their order, a rhetorical relation for each (elaboration, contrast, cause, comparison with history) and a hedge.
Decoding is constrained by a grammar generated per capsule from the ledger identifiers~\cite{ggerganov_llamacpp}, so the plan is always well-formed and cannot reference a unit or event that does not exist.
\emph{Ground:} deterministic code expands each plan item into a fact frame with values, names and units taken from $\mathcal L$.
\emph{Realize:} the SLM writes text from the fact frames, either with placeholders for numbers and names that are filled afterwards, or as free text.
\emph{Verify:} claims extracted from the text (Section~\ref{sec:metrics}) are checked against $\mathcal L$.
A failing sentence is regenerated at most twice and then replaced by the template realization of the same plan item.
The SLM therefore controls selection, ordering, aggregation and phrasing, but not values.

\paragraph{Hedging.}
The plan assigns hedges from fact confidence $p$: assertive for $p\ge0.9$, \emph{likely} for $0.7\le p<0.9$, \emph{possibly} for $0.4\le p<0.7$; facts with $p<0.4$ are mentioned only when asked about.
The verifier rejects a sentence whose hedge is stronger than its confidence band allows.
